\chapter{Cómo se fija una línea de ribera}

\label{cap:ribera}
La línea de ribera es el límite entre el dominio público del agua y la propiedad privada
del suelo. Determinarla es un acto administrativo: una comisión técnica la calcula, la
Autoridad de Aplicación la aprueba, se publica por edicto y desde entonces rige. De ese
acto depende cuánta superficie de un inmueble ribereño puede lotearse y cuánta no.

Este capítulo reconstruye el procedimiento, el programa que lo aplicó al departamento entre
2013 y 2015, y las dos determinaciones que se apartaron de la práctica general. Es la
pieza técnica del libro y conviene leerla antes que el capítulo siguiente, porque explica
qué sabía el Estado sobre estos ríos, desde cuándo, y qué de eso llegó a conocimiento
público.

\section{2025: la práctica continúa}

El grueso de las determinaciones que este capítulo analiza va de 2013 a 2015. Más adelante se
documenta una rectificación de junio de 2025 sobre este mismo departamento ---la Resolución 130/2025, que rectifica parcialmente una línea de 2012---; el boletín que
sigue es dos meses anterior y muestra que la práctica no se interrumpió en el medio.

\begin{ficha}{Boletín Oficial Nº 21.931, 14 de abril de 2025}
La Secretaría de Recursos Hídricos hace saber que por \textbf{Resolución Nº 000056/2025} del
20 de marzo de 2025 se aprobó la determinación de línea de ribera realizada por la Comisión
Técnica «en ambas márgenes del \textbf{Arroyo Vaqueros}, a la latitud del \textbf{Catastro
Nº 1122} del Dpto.\ La Caldera». Expediente 0090034-185051/2024-0.

«Cuyos \textbf{puntos de vinculación se encuentran a disposición de los interesados en la sede
de la Secretaría de Recursos Hídricos}.»

Autorizado a diligenciar: Ing.\ Juan Carlos Suárez. Firma la Dra.\ Silvia F.\ Santamaría\index{Santamaría, Silvia F.},
Jefa de Programa Jurídico.

\textbf{La misma resolución} aprueba una determinación sobre el \textbf{río Ancho, en el
departamento Cerrillos}, colindante a tres matrículas, en otro expediente iniciado en 2018.
\end{ficha}

\textbf{Primera, y es la que importa: no se publicaron las coordenadas.} El edicto remite los
puntos de vinculación a la sede del organismo. Es exactamente la fórmula de los dos actos de
2014 que este capítulo señala como excepción frente a las nueve que sí las publicaron en el mismo tramo.
\textbf{Once años después, la práctica que el libro registraba como minoritaria sigue en
uso.}

\textbf{Segunda: una sola resolución resuelve dos ríos de dos departamentos.} La 000056/2025
cubre el arroyo Vaqueros en La Caldera y el río Ancho en Cerrillos. Son expedientes distintos
--- uno iniciado en 2024, el otro en 2018 --- resueltos por el mismo acto el mismo día.

\textbf{Tercera: el trámite duró menos de un año en un caso y siete en el otro.} El expediente
del arroyo Vaqueros es de 2024 y se resolvió en marzo de 2025; el del río Ancho es de 2018.
El capítulo documenta que el trámite de Las Vertientes ---comisión en julio de 2013, determinación en abril de 2014--- duró ocho meses y medio, y que la determinación llegó sesenta y siete días después de que el plano la exigiera; la comisión del arroyo Quintín de 2013 recién se cerró en 2025.
\textbf{Los plazos de este procedimiento varían entre meses y más de una década, y el libro no
encontró criterio que lo explique.}

\textbf{Y cuarta: el edicto cita mal la norma.} El del arroyo Vaqueros remite al «Dcto.\ Nº
1898/02» y el del río Ancho, en el mismo boletín y bajo la misma resolución, al «Dcto.\ Nº
1989/02». El correcto es el segundo. Es un error de tipeo sin consecuencia jurídica, pero
conviene anotarlo porque \textbf{el edicto es la única vía por la que un tercero se entera de
que puede apelar}, y le indica el decreto equivocado.

\section{El diagnóstico del propio organismo}

Hay una pieza que dice, en boca del organismo que aplica la norma, casi todo lo que este
capítulo reconstruye desde afuera. Es una presentación del geólogo \textbf{Jorge G.\ Torres\index{Torres, Jorge G.}},
de la Secretaría de Recursos Hídricos de Salta, sobre línea de ribera y ordenamiento
territorial, publicada por el Consejo Hídrico Federal.

\subsection*{Primero, dos piezas del procedimiento que faltaban}

\textbf{La línea de ribera puede rectificarse.} El artículo 126 del Código de Aguas, que el
Decreto 1.989/02 reglamenta, termina con una facultad que este libro no había consignado: «La
Autoridad de Aplicación podrá \textbf{rectificar la línea de ribera cuando el cambio de las
circunstancias lo haga necesario}».

\textbf{Y se nombra el ``Certificado de No Inundabilidad''.} Lo emite el área de
Infraestructura Hídrica, y es «trámite necesario para loteos o urbanizaciones» cuyo límite sea
un río. Más adelante este capítulo obtiene el texto que lo funda --- el artículo 172 de la
reglamentación del Código de Aguas, que lo califica de \textbf{recaudo indispensable} --- y
esta presentación permite fecharlo como práctica corriente del organismo. \textbf{No aparece
nombrado en ninguno de los expedientes de loteo que el capítulo \ref{cap:loteo} analiza.}

La presentación agrega cómo se encadenan las dos cosas: «La línea de ribera definirá en casi
todos los casos los lugares y qué tipo de defensa».

\subsection*{Y después, el conflicto, dicho por el Estado}

\begin{ficha}{Torres, J.\ G., Secretaría de Recursos Hídricos de Salta --- Consejo Hídrico Federal}
«Por lo general, en urbanizaciones, las reglamentaciones del municipio ronda un \textbf{5\,\%
para espacios verdes}. \ldots\ En este punto también es en algunos caso conflictivo, ya que
\textbf{algunos emprendedores quieren la modificación de la Línea de Ribera, para poder
cumplir con el 5\,\% requerido por el municipio}.»

«Un caso reciente, y aun no resuelto es una urbanización privada, donde originalmente
\textbf{el espacio verde fue ocupado por la construcción de duplex}.»

Y en las conclusiones:

«La determinación de la línea de ribera de los ríos Arenales y \textbf{Vaqueros}, esta hecha
casi en su totalidad. \textbf{Aun así, los emprendedores inmobiliarios, solicitan la
modificación de sectores para ganar espacios y avanzar hacia el dominio fluvial.}»

«La \textbf{falta de coordinación, entre organismos provinciales y/o municipales}, se ve
reflejada en los constantes reclamos y problemas que se están generando.»
\end{ficha}

Cuatro cosas, y las cuatro son del Estado y no de este libro.

\textbf{Primera: el mecanismo del conflicto queda explicado.} El municipio exige un cinco por
ciento de espacio verde; la línea de ribera resta superficie utilizable; \textbf{y entonces se
pide correr la línea para que la cuenta cierre}. No es una sospecha de este trabajo: es la
descripción que hace el organismo que resuelve esos pedidos.

\textbf{Segunda: la línea de ribera de Vaqueros está determinada casi en su totalidad.} Eso
confirma desde adentro lo que el capítulo reconstruyó acto por acto, y agrega que el trabajo
está prácticamente completo sobre ese río.

\textbf{Tercera, y es la frase más fuerte: «avanzar hacia el dominio fluvial».} El organismo
describe una presión sostenida del mercado inmobiliario sobre el límite del dominio público.
\textbf{Este libro sostiene, en su primera tesis, que hay dos circuitos de suelo con
regulaciones distintas; la Secretaría de Recursos Hídricos describe uno de ellos empujando el
límite del otro.}

\textbf{Y cuarta: la falta de coordinación la declara el propio Estado.} No hace falta que este
trabajo la infiera del expediente: está en las conclusiones de una presentación oficial ante
el consejo federal del sector.

\begin{cautela}
La presentación no está fechada con precisión --- el archivo se publicó en 2022 y contiene
imágenes rotuladas 2019 y 2021 ---, no identifica el caso de la urbanización con los duplex, y
no dice cuántos pedidos de modificación de línea se presentaron ni cuántos prosperaron.

Este libro \textbf{no puede establecer si alguno de los emprendimientos que documenta pidió
correr su línea de ribera}. Sí puede afirmar tres cosas: que la práctica existe y el organismo
la describe como habitual y conflictiva; que \textbf{la norma la permite}, porque el artículo
126 autoriza a rectificar cuando cambien las circunstancias; y que \textbf{hay al menos una
rectificación fechada en este departamento}, que un apartado posterior documenta.
\end{cautela}

\section{Los aforos existieron, y están publicados}

Este libro repite, desde la declaración estatal de 1924 y a lo largo de un siglo, que nadie aforó estos ríos. Hay que corregirlo, y la corrección es peor que la afirmación original.

El Plan de Manejo de 2024 trae, en sus páginas 11 y 12, dos tablas con las estaciones de
medición de la cuenca. Las cinco de aforo son éstas:

\needspace{8\baselineskip}
{\sffamily\bfseries\footnotesize Estaciones de aforo de la cuenca --- Plan de Manejo, tabla 2\par}
\vspace{0.5em}
{\small
\setlength{\LTpre}{0.6em}\setlength{\LTpost}{1.2em}
\begin{longtable}{@{}p{3.9cm}p{2.4cm}p{1.4cm}p{2.0cm}p{2.2cm}@{}}
\toprule
\textbf{Estación} & \textbf{Río} & \textbf{Titular} & \textbf{Período} & \textbf{Área de aporte} \\
\midrule
\endhead
El Volcán & Las Nieves & SSRH & 1947--1961 & 300,55 km² \\
\textbf{El Angosto} & \textbf{Mojotoro} & SSRH & \textbf{1942--1986} & \textbf{828,34 km²} \\
San Alejo & San Alejo & SSRH & 1944--1955 & 38,10 km² \\
Santa Rufina & Santa Rufina & SSRH & 1944--1955 & 84,87 km² \\
Desembocadura al Nieves & Yacones & SSRH & 1947--1961 & 50,85 km² \\
\bottomrule
\end{longtable}
}

Y las cinco pluviométricas --- El Angosto, Las Nieves-Volcán, \textbf{Los Yacones}, Santa
Rufina y San Alejo --- están todas cerradas también. La última dejó de operar en 1990.
\textbf{El estudio de 2024 no dispone de una sola estación en operación dentro de la zona.}

Cuatro cosas se siguen de ahí.

\textbf{Primera: el dato existió, pero ninguna serie se acerca al requisito legal tanto como parece.} El Angosto, sobre el Mojotoro, \textbf{abarca cuarenta y cuatro años} ---de 1942 a 1986--- y ochocientos veintiocho kilómetros cuadrados de área de aporte. El artículo 8 del Decreto 1.989/02 pide series de cincuenta años. Por calendario, faltaron seis; por datos, bastante más, como muestra la serie descargada más abajo. Y es, de las cinco, la más alejada de la zona urbana de La Caldera.

\textbf{Segunda, y es la que corrige de raíz una impresión de este libro: las estaciones de aforo eran nacionales, no provinciales.} La columna de titularidad dice SSRH ---Subsecretaría de Recursos Hídricos de la Nación--- en las cinco de aforo de la tabla y en tres de las cinco pluviométricas. \textbf{La única de origen provincial en todo el corpus es Los
Yacones}, de la Dirección General de Aguas de Salta.

Eso reordena el reproche. Durante buena parte del siglo, \textbf{quien midió estos ríos fue el
Estado nacional}, y la provincia sostuvo una sola estación.

\textbf{Tercera: los registros están publicados.} La consultora no los recuperó de los
organismos titulares ni de archivo alguno: los tomó de dos compilaciones nacionales editadas
--- la \textit{Estadística Hidrológica de la República Argentina}, edición 2004, y el
\textit{Atlas Digital de los Recursos Hídricos Superficiales} del Instituto Nacional del
Agua.

\textbf{Los caudales de estos ríos están impresos en dos publicaciones oficiales desde hace veinte años.} No hay que buscarlos: hay que abrirlos.

Y se pueden descargar. El Sistema Nacional de Información Hídrica de la Subsecretaría de Recursos Hídricos conserva las cinco estaciones de la tabla dentro de la \textbf{Red Hidrológica Nacional}, entre sus estaciones inactivas, con su serie de caudal medio mensual.

\needspace{10\baselineskip}
{\sffamily\bfseries\footnotesize Caudales medios mensuales de las cinco estaciones --- SNIH, datos históricos, consultados en septiembre de 2026\par}
\vspace{0.5em}
{\small
\setlength{\LTpre}{0.6em}\setlength{\LTpost}{1.2em}
\begin{longtable}{@{}>{\raggedright\arraybackslash}p{3.2cm}>{\raggedright\arraybackslash}p{2.4cm}>{\raggedright\arraybackslash}p{1.7cm}>{\raggedright\arraybackslash}p{1.6cm}>{\raggedright\arraybackslash}p{2.5cm}@{}}
\toprule
\textbf{Estación (Nº SNIH)} & \textbf{Serie} & \textbf{Meses con dato} & \textbf{Caudal medio (m³/s)} & \textbf{Máximo mensual (m³/s)} \\
\midrule
\endhead
San Alejo (634) & sep.\ 1944 -- ago.\ 1955 & 132 de 132 & 1,89 & 16,15 (ene.\ 1947) \\
Santa Rufina (635) & sep.\ 1944 -- ago.\ 1955 & 132 de 132 & 2,71 & 16,98 (feb.\ 1952) \\
El Volcán, Las Nieves (623) & dic.\ 1947 -- ago.\ 1961 & 165 de 165 & 4,97 & 23,20 (feb.\ 1953) \\
Desembocadura al Nieves, Yacones (640) & dic.\ 1947 -- ago.\ 1961 & 165 de 165 & 1,24 & 6,80 (feb.\ 1959) \\
\textbf{El Angosto, Mojotoro (625)} & \textbf{sep.\ 1942 -- ago.\ 1986} & \textbf{484 de 528} & \textbf{14,70} & \textbf{121,5 (feb.\ 1974)} \\
\bottomrule
\end{longtable}
}

Las cuatro series cortas están completas y confirman los períodos de la tabla del Plan. La de El Angosto no, y eso corrige la cifra que este capítulo usaba. La lámina \ref{fig:caudales} muestra los huecos año por año y el régimen de las cinco estaciones.

\begin{figure}[htbp]
\centering
\includegraphics[width=\textwidth]{fig-caudales.png}
\caption[Caudales de la cuenca del río La Caldera, 1942--1986]{\textbf{Caudales de la cuenca, 1942--1986.} A: caudal medio de El Angosto por año hidrológico; en azul claro, los años con meses sin dato; las cruces marcan los dos años sin ningún registro. B: régimen mensual de las cinco estaciones, expresado como cociente entre el caudal de cada mes y el caudal medio del período. Elaboración propia con datos del SNIH, con las fechas corregidas.}\label{fig:caudales}
\end{figure}

\begin{ficha}{El Angosto, estación 625: qué hay dentro de los cuarenta y cuatro años}
\begin{itemize}[leftmargin=*,nosep]
\item \textbf{Faltan cuarenta y cuatro meses} de los quinientos veintiocho del período. Se concentran en dos tramos: 1967--1969, con veinticuatro meses sin dato, y 1975--1982, con veinte.
\item \textbf{Treinta y cinco de los cuarenta y cuatro años hidrológicos}, de septiembre a agosto, tienen los doce meses completos.
\item \textbf{Un defecto de la exportación, no de la serie.} Siete registros vienen fechados a las 23 horas del último día del mes anterior ---por ejemplo, el 28 de febrero de 1946--- y parecen, a primera vista, meses duplicados junto a meses faltantes. Corregida esa hora de desfase, cada valor cae en su mes y no queda ningún duplicado. El mismo defecto afecta un registro de San Alejo y otro de Santa Rufina. Los valores de este capítulo usan las fechas corregidas.
\item \textbf{Año más seco completo, 1948}, con 7,57~m³/s de media; \textbf{el más húmedo, 1974}, con 23,46, el mismo del máximo mensual de la serie.
\end{itemize}
\end{ficha}

Tres cosas se siguen de las planillas. La primera es que el dato no sólo existió: \textbf{está hoy a un clic, en una base pública del Estado nacional}, con fecha, variable y unidad. La segunda es que la serie más larga está más lejos del umbral legal de lo que su calendario sugiere: cuarenta y cuatro años de extensión, pero treinta y cinco completos. La tercera es que las series dicen algo sobre el río que ningún acto de este libro menciona. El régimen es estival y extremo: el caudal medio de febrero multiplica el de septiembre por nueve en El Volcán, por catorce en Santa Rufina, por diecisiete en San Alejo, por veintidós en El Angosto y por veintisiete en el Yacones. Y un solo mes llegó a promediar más de dieciséis metros cúbicos por segundo en cuencas de treinta y ocho y ochenta y cinco kilómetros cuadrados. Es el comportamiento que el Decreto 1.989/02 manda considerar, y que una determinación por método geomorfológico sólo puede inferir.

\textbf{Y cuarta: esa procedencia consta al pie de dos páginas de tablas, y en ningún otro
lado.} No figura en el capítulo de recopilación de antecedentes, que enumera nueve organismos
consultados sin decir qué aportó cada uno, ni en el de fuentes de información del sistema de
información geográfica, que lista doce fuentes y no menciona ninguna de las dos.

Es, en un documento técnico y en miniatura, exactamente lo que este libro llama la segunda
tesis: \textbf{el dato no está negado ni oculto; está en el lugar donde nadie lo busca.}

\begin{cautela}
Lo que este libro corrige es su propia afirmación. \textbf{No es cierto que nunca se hayan
aforado estos ríos}: se los aforó durante décadas, los registros se publicaron y siguen
disponibles.

Lo que sigue siendo cierto, y ahora se entiende mejor, es que el artículo 8 del Decreto
1.989/02 exige cincuenta años de serie y la mejor de estas estaciones abarca cuarenta y cuatro, con treinta y cinco completos, y cerró en 1986. \textbf{La determinación de línea de ribera se resuelve por método
geomorfológico no porque falte el dato, sino porque el dato que hay no alcanza el umbral que
la propia norma fijó} --- y no alcanza porque las estaciones cerraron.

Este libro descargó las cinco series, pero no verificó si alguna de las determinaciones que analiza las consultó. Ninguna las menciona. Tampoco procesó las series más allá de los valores que consigna la ficha: calcular caudales de crecida con períodos de retorno exige los caudales máximos instantáneos, no los medios mensuales, y eso es trabajo de un hidrólogo.
\end{cautela}

\subsection*{Y una sigue midiendo, treinta kilómetros aguas abajo}

La consulta del catálogo público de estaciones del Sistema de Información Hidrológica ---
mil novecientas treinta y siete estaciones de la Argentina --- permite precisar qué queda hoy
de esa red.

\textbf{De las cinco estaciones de aforo de la cuenca, ninguna figura en el catálogo activo.}
El Volcán, San Alejo, Santa Rufina y Desembocadura al Nieves no aparecen con ningún nombre ni
sobre ningún río. Tampoco hay estación alguna llamada La Caldera, Wierna, Vaqueros o Yacones.

\textbf{Pero el río Mojotoro sí tiene medición, y doble.}

\needspace{6\baselineskip}
{\small
\begin{longtable}{@{}p{3.6cm}p{2.4cm}p{2.6cm}p{3.4cm}@{}}
\toprule
\textbf{Estación} & \textbf{Tipo} & \textbf{Red} & \textbf{Registro} \\
\midrule
\endhead
Mojotoro -- Güemes & hidrológica & \mbox{Red Hidrológica} \mbox{Nacional, SSRH} & manual, base BDHI \\
Mojotoro -- Güemes & combinada & RHN -- SAT & \textbf{automática} \\
\bottomrule
\end{longtable}
}

Están a unos treinta y dos y treinta y seis kilómetros de La Caldera, ya en el departamento
General Güemes, aguas abajo de la confluencia. Una registra altura de forma manual y vuelca a
la base hidrológica nacional; la otra es \textbf{automática y forma parte del sistema de
alerta}.

Tres consecuencias, y la última corrige el reproche que este libro venía haciendo.

\textbf{Primera: la capacidad técnica existe y está desplegada en la región.} Dentro de un
radio de sesenta kilómetros hay \textbf{treinta y dos estaciones} de al menos cuatro redes
distintas --- la hidrológica nacional, el sistema de alerta, la del INTA y redes
meteorológicas ---, la mayoría automáticas. No falta equipamiento ni institución: falta
dentro de esta cuenca.

\textbf{Segunda: lo que se mide es el río después, no el río que inunda.} El Mojotoro en
Güemes recoge el aporte de toda la cuenca alta, incluido el La Caldera. Sirve para saber qué
bajó; no sirve para avisarle a La Caldera que algo está bajando, porque para cuando el agua
llega a Güemes ya pasó por el pueblo.

Es exactamente la distinción que el Plan de Manejo hace al proponer su sistema de alerta:
el punto útil está \textbf{cinco kilómetros aguas arriba} del pueblo, porque de ahí la
crecida tarda entre cuarenta y cinco y sesenta minutos en llegar. \textbf{Medir treinta
kilómetros aguas abajo es medir demasiado tarde.}

\textbf{Y tercera: el río de este departamento nunca dejó de estar medido del todo.} Lo que se
abandonó fue la medición \textit{dentro} de la cuenca. La estación de El Angosto, sobre el
mismo Mojotoro, funcionó de 1942 a 1986; hoy hay dos estaciones sobre ese río, más abajo, una
de ellas automática.

\begin{cautela}
El catálogo consultado publica la existencia y la ubicación de cada estación, no sus series.
Este libro \textbf{no consultó los registros} de las dos estaciones del Mojotoro, ni sabe
desde cuándo operan ni qué continuidad tienen.

Tampoco puede afirmar que esos datos sirvan para determinar una línea de ribera en La Caldera:
el área de aporte es distinta y el punto de medición está fuera del departamento. Lo que
registra es más modesto y más incómodo: \textbf{el Estado sigue midiendo este sistema hídrico,
sólo que en el tramo donde el agua ya no es problema de nadie de este departamento.}
\end{cautela}

\section{El procedimiento, leído}

Los veinticinco actos que este capítulo analiza invocan todos la misma norma ---varios edictos, con el número mal escrito: «1898/02»---: el \textbf{Decreto
Nº 1.989/02}, que aprueba la Resolución 070/02 de la Agencia de Recursos Hídricos y reglamenta
el artículo 126 del Código de Aguas. Hasta ahora este libro la conocía sólo por lo que decían
los actos que la citaban. Conviene leerla, porque contesta preguntas que el capítulo venía
dejando abiertas.

\textbf{Y conviene empezar por cómo se publicó}, porque la forma de la publicación dice tanto
como el contenido.

\begin{ficha}{Boletín Oficial Nº 16.515, 11 de noviembre de 2002, pág.\ 4.773}
Sección Administrativa, Decretos. La entrada del sumario reza: «\textsc{m.p.e.\ n\textup{º} 1.989 del
28/10/02 --- Código de Aguas de la Provincia de Salta, Ley N\textup{º} 7.017: art.\ 126, determinación
Líneas de Rivera}».

El cuerpo de la sección trae, bajo el título \textbf{DECRETO Nº 1.989}, un solo párrafo:

\begin{quote}
«El texto del Decreto Nº 1989/02 que reglamenta el art.\ 126 de la Ley Nº 7017 --- Código de
Aguas de Salta, \textbf{se publica como Separata de este Boletín Oficial}.»
\end{quote}

Y debajo, la fecha: \textbf{Salta, 28 de octubre de 2002}. Nada más.
\end{ficha}

\begin{cautela}
Dos cosas.

\textbf{Primera, las fechas y la edición.} El decreto es del \textbf{28 de octubre de 2002}; el 11 de noviembre es la fecha de la edición que lo publica, el \textbf{Boletín Oficial Nº 16.515}, página 4.773.

\textbf{Y segunda:} la edición del Boletín \textbf{no contiene el decreto}. Contiene una línea
que dice dónde está. El acto se publica; el texto que permitiría cumplirlo, verificarlo o
discutirlo va a un impreso aparte, que se distribuye aparte y que no forma parte del archivo
digitalizado de la edición.

\textbf{Es la estructura que el capítulo \ref{cap:opacidad} describe, en su forma más concentrada y once años antes de la serie de determinaciones que este capítulo reconstruye.} Nada de esto
es irregular: la separata es un instrumento previsto y legítimo. Lo que se registra es el
efecto. La norma que rige los veinticinco actos de línea de ribera de este departamento --- la que fija quién integra la comisión, con qué series de aforo, con qué
publicación y a costa de quién --- \textbf{no está donde dice el acto que la publica}.
\end{cautela}

\begin{ficha}{Decreto Nº 1.989/02 --- Anexo I, procedimiento en doce artículos}
\textbf{Art.\ 1º} El trámite se inicia \textbf{de oficio o a pedido de parte}.
\textbf{Art.\ 3º} La Comisión Técnica se integra con un \textbf{ingeniero hidráulico o civil,
un geólogo y un agrimensor}.
\textbf{Art.\ 4º} Notificación y \textbf{publicación por dos días} en el Boletín Oficial y en
el diario de mayor difusión.
\textbf{Art.\ 5º} Quince días para que las partes presenten información.
\textbf{Art.\ 7º} La delimitación se hace por el caudal de la \textbf{máxima crecida
ordinaria}, considerando la \textbf{migración lateral por erosión natural}.
\textbf{Art.\ 8º} Prelación de fuentes: \textbf{aforos sistemáticos con series no menores de
cincuenta años}; en su defecto, registros fotográficos y satelitales; y recién después,
modelos matemáticos de cuenca.
\textbf{Art.\ 9º} Trabajo de campo: \textbf{amojonar} a izquierda y derecha los lugares de las
más altas aguas por donde naturalmente el río ha pasado.
\textbf{Art.\ 10} La Comisión propone, la Autoridad notifica, los interesados tienen diez días
para observar, la Autoridad resuelve en diez días hábiles, \textbf{se registra en el libro de
Catastros y se comunica a la Dirección de Inmuebles, previa publicación por dos días}. El acto
es apelable ante el \textbf{Tribunal de Aguas}.
\textbf{Art.\ 11} \textbf{Los costos los absorbe quien solicita} la determinación.
\end{ficha}

Cinco cosas quedan establecidas, y tres de ellas cambian lo que este capítulo podía afirmar.

\textbf{Primera: la publicación es obligatoria y es parte del acto.} El artículo 10 no dice
que se publique el resultado: dice que la resolución «se registra en el libro de Catastros y
se comunica a la Dirección de Inmuebles, \textbf{previa publicación por dos días}». La
publicación es condición del registro, no un trámite posterior.

\begin{cautela}
El decreto \textbf{no dice expresamente que deban publicarse las coordenadas}. Dice que debe
publicarse la resolución, y que el trabajo de campo consiste en amojonar los puntos de
ribera. Este libro \textbf{no puede afirmar}, entonces, que las tres resoluciones que
remitieron sus coordenadas a la sede del organismo hayan incumplido el decreto.

Lo que sí puede afirmar es que \textbf{doce de las quince determinaciones de
2013--2015 publicaron sus coordenadas y tres no}, bajo la misma norma y en el mismo período, y las tres con el mismo asesor legal firmante. Y que las dos de esas tres que se refieren a un inmueble recaen sobre suelo que este libro
encontró en proceso de loteo ---condición que comparten, sin embargo, otras dos que sí las publicaron---. La norma admite las dos prácticas; lo que no explica es por qué
se eligió una u otra en cada caso, y eso es exactamente lo que la segunda tesis sostiene.
\end{cautela}

\textbf{Segunda: los aforos están en primer lugar de prelación, y con series de cincuenta
años.} El artículo 8 los pone antes que cualquier otro método. El capítulo \ref{cap:siglo} documenta que en 1924 el propio Estado dejó constancia escrita de no haber aforado ningún río de la provincia.
\textbf{La norma exige series de cincuenta años, y la más larga de las cinco estaciones de la cuenca ---El Angosto, 1942 a 1986--- abarca cuarenta y cuatro, de los que treinta y cinco tienen los doce meses medidos; las otras cuatro cerraron entre 1955 y 1961 con once a catorce años de registro completo.}

\textbf{Tercera: el decreto contempla que el río se mueva.} El artículo 7 manda considerar la
«migración lateral por erosión natural». Es la misma constatación que en 1940 obligó a
suspender una obra porque «el río desvió su curso», ahora incorporada al procedimiento.

\textbf{Cuarta: el acto es apelable ante el Tribunal de Aguas.} Existe una vía de revisión
específica. Este libro no encontró ningún caso del departamento que la haya usado.

\textbf{Y quinta: los costos los paga quien solicita.} Cuando el trámite lo inicia un
propietario que quiere lotear, es él quien financia la determinación de hasta dónde puede
lotear. La norma lo dice sin rodeos, y no es una anomalía: es el régimen. Pero conviene
leerlo junto al artículo 1, que permite iniciar \textbf{de oficio}: \textbf{el Estado puede
determinar la línea de ribera sin que nadie se lo pida y a su propio costo}, y este libro no
encontró constancia de que lo haya hecho en este departamento.

\section{La determinación es anterior al aluvión}

Hay un antecedente administrativo anterior a todo lo demás, y es el más incómodo del expediente.

\begin{ficha}{Boletín Oficial Nº 19.467 (2015) --- sección Determinación de Línea de Ribera}
\begin{itemize}[leftmargin=*,nosep]
\item \textbf{Orden de publicación:} 100045883, página 407.
\item \textbf{Organismo:} Secretaría de Recursos Hídricos.
\item \textbf{Expediente:} Nº \textbf{34-167.127/13}.
\item \textbf{Objeto:} \textbf{Arroyos El Durazno y Guaranguay}, inmueble \textbf{catastro Nº 4.366}, departamento La Caldera.
\end{itemize}
\end{ficha}

El cuerpo del edicto fue obtenido. Su texto establece lo siguiente:

\begin{ficha}{Edicto --- Resolución Nº 383 del 18/12/2014}
La Secretaría de Recursos Hídricos hace saber que \textbf{por Resolución Nº 383 del 18 de diciembre de 2014 se aprobó la determinación de la línea de ribera} realizada por la Comisión Técnica \textbf{en ambas márgenes de los arroyos El Durazno y Guaranguay}, en la zona del inmueble identificado con la matrícula catastral Nº 4.366 del departamento La Caldera.

Las coordenadas Posgar de las líneas de ribera determinadas \textbf{``se encuentran a disposición de los interesados en la sede de la Secretaría''}.

El acto es apelable ante el \textbf{Tribunal de Aguas} --- Decreto Nº 1898/02, art.\ 10 --- dentro de los diez días hábiles administrativos contados desde la última publicación, presentando el recurso en Avenida Bolivia 4650. Publicado los días 21 y 22 de enero de 2015 en los \textbf{Boletines Oficiales Nº 19.466 y 19.467}, orden de publicación 100045883. Salta, 20 de enero de 2015. Suscribe el edicto el Dr.\ Matías J.\ Brogin\index{Brogin, Matías J.}, asesor legal.
\end{ficha}

\begin{cautela}
No es un trámite iniciado ni pendiente: \textbf{es una determinación aprobada}. El Estado provincial fijó, por resolución firme, dónde estaba el límite del cauce en ambas márgenes de los dos arroyos.
\end{cautela}

Y la fecha es exacta hasta el punto de resultar incómoda de escribir. La resolución que aprueba la línea de ribera es del \textbf{18 de diciembre de 2014}. El primer aluvión que destruyó los loteos fue el \textbf{21 de diciembre de 2018}.

\textbf{Cuatro años y tres días.}

\begin{cautela}
El edicto agrega, además, uno de los casos del patrón que el capítulo \ref{cap:opacidad} documenta.

Lo que se publica es que la línea existe. \textbf{Dónde está} --- las coordenadas Posgar que la definen --- no se publica: queda a disposición de los interesados en la sede de la Secretaría, en avenida Bolivia 4650, primer piso, ciudad de Salta.

De modo que el acto por el cual el Estado determina hasta dónde llega el cauce de un arroyo, y por lo tanto dónde no se puede construir, se comunica sin el dato que lo hace aplicable. Un comprador de lote en El Durazno podía leer ese edicto en enero de 2015 y no saber si su parcela estaba dentro o fuera de la línea.

Corresponde señalar que el mismo Boletín publica, para otro expediente de la misma sección, la tabla completa de coordenadas. Es decir: la práctica de publicarlas existía. En este caso no se aplicó.
\end{cautela}

\begin{cautela}
Subsiste el límite ya señalado y ahora es más preciso. El inmueble es la \textbf{matrícula catastral Nº 4.366}; el loteo El Durazno se asienta sobre la \textbf{matrícula Nº 6.558}. Este libro no establece la relación entre ambas.

Lo que puede afirmarse sin ese dato es: la línea de ribera de ambas márgenes de los arroyos El Durazno y Guaranguay estaba determinada por resolución desde diciembre de 2014, y el arroyo Guaranguay es aquél cuya crecida la sentencia identifica como causa del aluvión, aquél que en enero de 2026 volvió a aislar a los barrios del sur, y aquél sobre el cual se construyó el badén de ciento veinte millones de pesos.
\end{cautela}

\section{El procedimiento, y la fuente que lo describe}

El material que este capítulo venía usando con reservas quedó identificado. Se trata de una presentación institucional titulada \textit{Línea de Ribera y Ordenamiento Territorial}, cuyo autor es el \textbf{geólogo Jorge G.\ Torres\index{Torres, Jorge G.}, de la Secretaría de Recursos Hídricos de Salta}, publicada por el \textbf{Consejo Hídrico Federal (COHIFE)}. La reserva se levanta: es un documento oficial, de autor identificable, y su autor figura en diciembre de 2025 como \textbf{Director General de Cuencas y Política Hídrica} de la Subsecretaría --- es decir, el funcionario a cargo del área donde vive el Plan de Manejo de esta causa.

\begin{cautela}
La presentación agrega un dato que reordena el capítulo entero.

\begin{quote}
``La Secretaría de Recursos Hídricos, en el área de Infraestructura Hídrica, es la que emite el \textbf{Certificado de No Inundabilidad}. Trámite necesario para loteos o urbanizaciones, que algunos de los límites sea un río, se exige en primera medida la determinación de la línea de ribera, además de otros requisitos.''
\end{quote}

El loteo El Durazno está atravesado por los arroyos El Durazno y Guaranguay, cuya línea de ribera se determinó en 2014. El organismo que la determinó es el mismo que emite el Certificado de No Inundabilidad, y ese organismo intervino \textbf{a fojas 59} del expediente de aprobación del plano.

\textbf{La pregunta del capítulo deja de ser abierta y pasa a tener una única forma: ¿se emitió Certificado de No Inundabilidad para El Durazno, y en qué términos?} Si la foja 59 lo contiene, hay que leerlo. Si no lo contiene, hay que preguntar por qué un fraccionamiento limitado por dos cursos de agua se aprobó sin él. Este libro no sabe la respuesta y no la insinúa; señala que la pregunta ahora es contestable con una sola foja.
\end{cautela}

Y un segundo dato, sobre el visado. La presentación describe que, obtenida la resolución, ``el interesado puede realizar el plano y la Secretaría de Recursos Hídricos visa el mismo'', y enumera los requisitos: que los puntos de vinculación estén graficados, que las planillas lleven las coordenadas, y \textbf{que se mencione el número de Resolución}.

Eso vuelve verificable algo que hasta aquí no lo era. Si el plano 1.164 fue visado por Recursos Hídricos, debe mencionar en su carátula el número de la resolución de línea de ribera. \textbf{Basta mirar el plano para saber si la línea de 2014 fue incorporada al fraccionamiento de 2021.} Es una comprobación de un minuto, sobre un documento que ya está identificado.

\section{Lo que la Provincia ya hacía en 2012 y 2013}

Tres edictos anteriores a la serie muestran que las formas que después variaron ya estaban todas disponibles.

\begin{ficha}{Resolución Nº 94/12 --- B.O.\ Nº 19.130 y 19.131, 22 y 23 de agosto de 2013}
Expediente 0050034-71745/2011. Dictada el 2 de mayo de 2012; publicada quince meses después, con la transcripción literal de su parte dispositiva. Tiene por conformada la comisión técnica y determina la línea de ribera del \textbf{río Castellanos}, ambas márgenes, en la \textbf{matrícula 4027}, «cuya gráfica ilustra plano que forma parte del informe», con dos vértices por margen en coordenadas publicadas.

\textbf{Art.\ 3º:} advierte al propietario que, «\textbf{dada la peligrosidad del Río Castellanos}, que en oportunidad de crecidas extraordinarias podría rebasar eventualmente la línea de ribera demarcada, \textbf{deberá evitar la construcción de viviendas en la proximidad de la ribera}».

\textbf{Art.\ 4º:} la publicación es previa al registro de la línea «en el Libro de Catastros de la Secretaría de Recursos Hídricos» (art.\ 10 de la Resolución 070/02).

\textbf{Art.\ 5º:} comunica la resolución al Ministerio, al propietario, al ingeniero en construcciones Pablo Vicente Balverdi y \textbf{a la Dirección General de Inmuebles}. Firman el secretario Fuertes y la jefa del Programa Fiscalización y Control, Silvia F.\ Santamaría.
\end{ficha}

\begin{ficha}{Resolución Nº 155 del 05/06/2013 --- B.O.\ Nº 19.093, 28 de junio y 1 de julio de 2013}
Expediente 34-132.589/12. Determina la línea de ribera del \textbf{arroyo Los Nogales} en la \textbf{matrícula 1905}, sección A, fracción 17b ---el inmueble de la futura urbanización Cerro de Buena Vista---. «Las coordenadas de vinculación determinadas \textbf{se encuentran a disposición de los terceros interesados por mesa de entradas}.» Cita el artículo 126 del Código de Aguas y el Decreto 1989/02 con su número correcto. Firma Brogin.
\end{ficha}

\begin{ficha}{Resolución Nº 201 del 01/07/2013 --- B.O.\ Nº 19.109 y 19.110, 23 y 24 de julio de 2013}
Expediente 34-214.573/12. Determina la línea de ribera del \textbf{río Vaqueros}, ambas márgenes, en el inmueble de «matrícula catastral \textbf{S/Nº}», sección B, localidad de Vaqueros. \textbf{Publica dieciocho coordenadas Posgar}: trece en la margen izquierda y cinco en la derecha, a lo largo de unos trescientos sesenta metros. Cita el «Dcto.\ 1898/02». Firma Brogin.
\end{ficha}

\begin{cautela}
Tres observaciones, de peso desigual.

\textbf{La primera es la que más pesa.} En mayo de 2012 la Secretaría de Recursos Hídricos no se limitó a fijar una línea: advirtió por escrito al propietario que no construyera viviendas cerca de la ribera de un río peligroso, aunque la línea quedara demarcada, y comunicó la resolución a Inmuebles, el organismo que aprueba los planos. El instrumento para prevenir, en el mismo acto que fija la ribera, existía y se usaba dos años antes de la determinación de El Durazno. Este libro no encontró una advertencia equivalente en ningún acto sobre los arroyos El Durazno y Guaranguay; el edicto de la Resolución 383/14 no la trae, aunque, como se dijo, un edicto es sólo la síntesis de la resolución.

\textbf{La segunda es de forma.} Los tres edictos, del mismo año de publicación y con el mismo asesor en dos de ellos, muestran las tres variantes que la serie posterior alterna: coordenadas publicadas (201/13), coordenadas en mesa de entradas (155/13) y parte dispositiva transcripta con advertencias (94/12). También muestran que la cita equivocada del decreto convivía, en 2013, con la correcta. Una curiosidad de la 94/12: sus dos vértices de la margen izquierda distan entre sí cinco metros; los de la derecha, casi doscientos.

\textbf{La tercera es de lugar.} La línea del río Vaqueros de 2013 termina, en uno de sus extremos, a trece metros del primer punto de la ampliación que la Resolución 144/17 agregó en 2017 a la línea de la matrícula 40. Leídas juntas, las Resoluciones 201/13, 414/16 y 144/17 fijan más de seiscientos metros continuos del mismo río, con un ancho entre márgenes de cien a ciento cuarenta metros, en tres actos y cuatro años, para inmuebles vecinos. El de 2013 no tenía número de matrícula. Y ese tramo del río es límite departamental: según la cartografía de radios del Censo 2022, el punto medio entre ambas márgenes cae a seis metros (2013) y a trece (2016) del límite de La Caldera, del lado del departamento Capital. Las líneas de las dos orillas se fijaron en expedientes de La Caldera, aunque la margen derecha está, para el INDEC, en otro departamento. \pendiente{identificación del inmueble «s/Nº» de la sección B de Vaqueros y su relación con la matrícula 40}
\end{cautela}

\section{Dos años, veinticinco actos}

La determinación de diciembre de 2014 sobre El Durazno y Guaranguay no fue un acto aislado. Entre octubre de 2013 y diciembre de 2015 la Secretaría de Recursos Hídricos dictó veinticinco actos de comisión o de determinación sobre la cuenca del río La Caldera y sus afluentes; la tabla agrega la Resolución 023/14, que designa y determina a la vez.

\textbf{Conviene fijar la base del recuento, porque este capítulo la usa de dos maneras.} La tabla tiene veintiséis filas: \textbf{diez comisiones}, \textbf{quince determinaciones} ---incluida la 023/14, que la serie cuenta aparte por designar y determinar en un solo acto--- y una fila del índice del Boletín cuyo acto no se identificó. Cuando este libro dice «veinticinco actos» excluye a la 023/14; cuando dice «quince determinaciones» la incluye. Son dos recuentos distintos de la misma tabla y ninguno contiene al otro.

\needspace{6\baselineskip}
{\sffamily\bfseries\footnotesize Actos de línea de ribera en el departamento, octubre de 2013 -- diciembre de 2015 \textit{(por fecha de publicación; la columna consigna la fecha de dictado, y por eso la primera fila es de julio de 2013. Dos filas rompen ese orden a propósito: la Resolución 428/15, cuyo edicto salió en febrero de 2016, y la Resolución 383/14, que se deja al final por ser el acto que este capítulo persigue)}\par}
\vspace{0.5em}
{\small
\setlength{\LTpre}{0.6em}\setlength{\LTpost}{1.2em}
\begin{longtable}{@{}p{2.2cm}p{2.9cm}p{6.8cm}@{}}
\toprule
\textbf{Acto} & \textbf{Expediente} & \textbf{Objeto} \\
\midrule
\endhead
Res.\ 208/13 (22/07/13) & 34-240.329/12 & \textbf{Comisión}: arroyo Quintín o Pacará, catastro 3989, fracción 78. Integran el ing.\ civil César Soto, el geólogo Marcelo Ramón Olañeta y el técnico Félix Márquez. B.O.\ 20 y 21/11/2013 \\
Res.\ 292/13 (10/10/13) & 34-160.869/13 & \textbf{Comisión}: arroyo Quintín, catastro 3756, fracción 48. B.O.\ 27 y 28/03/2014. La determinación llega en 2025 (Res.\ 142/2025) \\
Res.\ 348/13 (19/11/13) & --- & \textbf{Determinación}: río La Caldera, matrícula 1590, La Calderilla. \textit{Con coordenadas y zona inundable publicadas} \\
Res.\ 395/13 (26/12/13) & 34-167.127/13 & \textbf{Comisión}: río La Caldera y arroyo del Manzano, catastro 4366 (el aviso dice «1366»), Finca La Caldera o Getsemaní \\
Res.\ 6/14 (13/01/14) & 34-179.209/12 & \textbf{Determinación}: río La Caldera, matrícula 3593, fracción 17. \textit{Con coordenadas publicadas} \\
Res.\ 7/14 (13/01/14) & 34-19.384/12 & \textbf{Determinación}: río La Caldera, matrícula 3584, La Calderilla. \textit{Con coordenadas publicadas} \\
Res.\ 8/14 (13/01/14) & 34-207.589/12 & \textbf{Comisión}: quebrada de la matrícula 4191, Fracción Finca Vaqueros \\
Res.\ 023/14 (06/02/14) & 90.034-96.837/13 & Designa al geólogo Jorge G.\ Torres\index{Torres, Jorge G.} y \textbf{determina} la línea de ribera del río La Caldera a lo largo del \textit{pueblo}; coordenadas en el anexo vinculado \\
Res.\ 80/14 (07/04/14) & 34-240.329/12 & \textbf{Determinación}: arroyo Quintín o Pacará, matrícula 3989, fracción 78, Finca Vaqueros --- \textit{el inmueble del Club de Campo Las Vertientes}. \textit{Con coordenadas publicadas} \\
Res.\ 96/14 (29/04/14) & 34-61.721/14 & \textbf{Comisión}: río Vaqueros, catastros 1905 sec.\ A frac.\ 17-b y 3740 frac.\ 47 \\
Res.\ 101/14 (29/04/14) & 34-207.589/12 & \textbf{Determinación}: arroyo Los Nogales Norte o El Cajón, matrícula 4191. \textit{Con coordenadas publicadas} \\
Res.\ 182/14 (03/07/14) & 34-160.829/13 & \textbf{Determinación}: río La Caldera, matrícula 2882, fracciones 1 y 3, La Calderilla. \textit{Con coordenadas publicadas} \\
\textbf{Res.\ 185/14} (11/07/14) & \textbf{34-167.127/13} & \textbf{Determinación}: río La Caldera y arroyo El Manzano, \textbf{matrícula 4366}, Finca La Caldera o Getsemaní. \textit{Con coordenadas publicadas} \\
Res.\ 142/14 (27/05/14) & 125-19.223/09 & \textbf{Determinación}: ríos Mojotoro y La Caldera, matrícula 3616 ---el inmueble del loteo La Misión Etapa III---. \textit{Sin coordenadas publicadas} \\
Res.\ 209/14 (29/07/14) & 34-61.721/14 & \textbf{Determinación}: río Vaqueros, matrículas 1905 y 3740. \textit{Con coordenadas publicadas} \\
Res.\ 74/15 (10/03/15) & 34-259.760/14 & \textbf{Comisión}: río La Caldera, catastros 3721 y 3722, Fracción Finca Wierna \\
Res.\ 197/15 (05/06/15) & 34-284.285/14 & \textbf{Comisión}: río La Caldera, catastro 4049. Integran el \textbf{geólogo Jorge Torres} y el ing.\ civil Carlos Cerezo. Método geológico-morfológico; el índice del Boletín titula el aviso con Aguas del Norte \\
Res.\ 198/15 (05/06/15) & 34-145.043/15 & \textbf{Comisión}: río Vaqueros, catastros 156, 23, 1695 y 406 \\
--- (jul.\ 2015) & 34-485.548/14 & Arroyo Isasmendi, Finca Wierna, catastro 2020 (sólo el título del índice) \\
Res.\ 428/15 (03/12/15) & 34-\textbf{1}85.548/14 & \textbf{Determinación}: \textbf{río La Caldera}, matrícula 2020, Fracción A, Finca Wierna; el texto dice «Dpto.\ La Candelaria». \textit{Con coordenadas publicadas} (trece puntos en la margen izquierda, cuatro en la derecha). Edicto del 22/12/2015, B.O.\ 10 y 11/02/2016 \\
Res.\ 273/15 (13/08/15) & 34-125.733/15 & \textbf{Comisión}: río Trancas, catastro 4445, Fracción 195, Finca Vaqueros. Integran Forns, el geólogo Teodoro Chafatinos y el ing.\ civil Lucas Vidal. La determinación resultante (Res.\ 14/16) nombra el río Wierna \\
Res.\ 293/15 (31/08/15) & 34-259.760/14 & \textbf{Determinación}: río La Caldera y Cañada Oeste, matrículas 3721 y 3722, Finca Wierna. \textit{Con coordenadas publicadas}. Edicto del 08/09/2015, B.O.\ Nº 19.666, 18 y 19/11/2015 \\
Res.\ 322/15 (25/09/15) & 34-145.043/15 & \textbf{Determinación}: río Vaqueros, matrículas 156, 23, 1695 y 406. \textit{Con coordenadas publicadas}. El edicto cita el expediente como 34-145.043/\textbf{14}; el de la comisión 198/15, como /15. Edicto del 26/10/2015, B.O.\ Nº 19.675, 2 y 3/12/2015 \\
Res.\ 319/15 (16/09/15) & 34-171.172/15 & \textbf{Comisión}: río Wierna, matrículas 4097, 4953, 391, 389, 1268, 3718 y 3438. Integran el geólogo Jorge Torres y el ing.\ Carlos Cerezo. Edicto del 26/10/2015, publicado recién el 21 y 22/12/2015 \\
\textit{Res.\ 467/15} (04/12/15) & 34-30.704/15 & \textit{Determinación}: río Lesser, tramo entre el puente de la R.P.\ 115 y su confluencia con el río Castellanos. \textit{Sin coordenadas publicadas} \\
Res.\ 383/14 (18/12/14) & 34-167.127/13 & \textbf{Determinación}: arroyos El Durazno y Guaranguay, matrícula 4366. \textit{Sin coordenadas publicadas} \\
\bottomrule
\end{longtable}
}

Firmados por el Dr.\ Matías J.\ Brogin\index{Brogin, Matías J.}, Asesor Legal, salvo el de la Resolución 023/14, que firma el secretario Alfredo Fuertes. Los que aprueban comisiones invocan el \textbf{Decreto 1989/02, art.\ 8 ap.\ 2 y 3}, prohíben a los ribereños modificar las márgenes mientras dure la operación y abren quince días hábiles para que terceros aporten estudios.

Hay cuatro cuestiones que esta serie resuelve.

\textbf{Primera: el edicto inicial existe, y este capítulo lo tiene.} Se venía pidiendo el acto que aprueba la comisión técnica, con nombres y método, previo a la determinación. Las Resoluciones 395/13 y 8/14 son exactamente eso, publicadas con la estructura completa: integrantes, inmueble, prohibición cautelar y plazo de quince días. La pregunta ya no es si ese acto se publica: se publica. La pregunta es si se publicó \textbf{para los arroyos El Durazno y Guaranguay}, y sigue sin aparecer.

\textbf{Segunda: las comisiones no tienen composición fija, y hay un nombre que se repite.} La 292/13 nombra al ingeniero civil \textbf{Antonio Alejandro Forns\index{Forns, Antonio Alejandro}} y al geólogo Teodoro Chafatinos. La 395/13 nombra al mismo Forns y al geólogo Mario Raúl Orte. La 208/13 nombra tres: el ingeniero civil César Soto, el geólogo Marcelo Ramón Olañeta y el técnico Félix Márquez. La 8/14 nombra tres: el ingeniero civil César Manuel Soto, el geólogo Jorge David Afranllie y el técnico universitario en topología Félix Márquez. La 96/14 nombra al ingeniero Wanny Caramella y a la geóloga Marcia Jesica Faviana. La 273/15, dos años después de la 292/13, vuelve a reunir a Forns y a Chafatinos, con un tercer integrante. La del río Bermejo, en 2019, nombra dos.

El artículo 3º del Anexo I prevé ingeniero hidráulico, geólogo y agrimensor. En la práctica la integración varía entre dos y tres miembros, la especialidad del ingeniero varía, y \textbf{en ninguna de las ocho comisiones cuya integración consta aparece un agrimensor}. Entre todas las comisiones que este capítulo leyó después, de 2011 a 2019, sólo una ---la del río Vaqueros de 2016--- incluye un ingeniero hidráulico. Es un dato sobre la aplicación de la norma, no sobre su texto, y merece preguntarse a la autoridad de aplicación.

\textbf{Tercera: el departamento estaba siendo relevado sistemáticamente.} Veinticinco actos en veintiséis meses, sobre más de veinticinco inmuebles y trece cursos distintos --- los ríos La Caldera, Vaqueros, Mojotoro, Wierna, Lesser y Trancas, la Cañada Oeste, y los arroyos Quintín o Pacará, Los Nogales Norte o El Cajón, El Manzano, Isasmendi, El Durazno y Guaranguay. La línea de ribera no se fijaba de manera ocasional: \textbf{había un programa en curso}, y el departamento entero estuvo bajo relevamiento hídrico entre 2013 y 2015, con los ciclos de comisión y determinación completándose uno tras otro.

Eso vuelve muy difícil sostener que la determinación sobre los arroyos del loteo El Durazno fuera un hecho excepcional, aislado o desconocido para la administración. Y agrega una pregunta que el capítulo no tenía: \textbf{¿por qué se relevó tanto y tan seguido justo entonces?} Un programa de esa intensidad suele responder a una decisión, a un financiamiento o a un conflicto previo. Cuál de los tres es materia del expediente administrativo del programa, no de los edictos.

\begin{cautela}
\textbf{Y hay un cuarto hallazgo, que es el que más pesa: el ciclo se completa y se publica entero.}

La Resolución 8/14 de enero aprobó la comisión técnica para ``la quebrada que atraviesa la matrícula 4191'', expediente 34-207.589/12. La Resolución 101/14, del 29 de abril, aprueba la determinación resultante sobre esa misma matrícula y ese mismo expediente, y el edicto \textbf{publica las coordenadas completas}. Entre una y otra pasaron tres meses y medio. La quebrada, además, deja de ser anónima: se llama \textbf{arroyo Los Nogales Norte o El Cajón}.

El mismo ciclo se ve en el arroyo Quintín o Pacará: la Resolución 208/13, de julio de 2013, designa la comisión sobre el catastro 3989, y la 80/14, de abril de 2014, aprueba la determinación con coordenadas, ocho meses y medio después, en el mismo expediente. No todos los ciclos se cerraron tan rápido: la comisión de la Resolución 292/13 sobre el catastro 3756, publicada en marzo de 2014, recién tuvo su determinación en junio de 2025.

Con eso puede hacerse el recuento que este capítulo necesitaba. De las \textbf{once determinaciones} de línea de ribera publicadas en el departamento entre noviembre de 2013 y diciembre de 2014, más una rectificación:
\end{cautela}

{\small
\setlength{\LTpre}{0.6em}\setlength{\LTpost}{1.2em}
\begin{longtable}{@{}p{2.1cm}p{6.6cm}p{3.2cm}@{}}
\toprule
\textbf{Resolución} & \textbf{Curso e inmueble} & \textbf{Coordenadas} \\
\midrule
\endhead
348/13 & Río La Caldera, matrícula 1590 & \textbf{Publicadas} \\
6/14 & Río La Caldera, matrícula 3593 & \textbf{Publicadas} \\
7/14 & Río La Caldera, matrícula 3584 & \textbf{Publicadas} \\
023/14 & Río La Caldera, a lo largo del pueblo & \textbf{Publicadas en el anexo} \\
80/14 & Arroyo Quintín, matrícula 3989 & \textbf{Publicadas} \\
101/14 & Arroyo Los Nogales Norte, matrícula 4191 & \textbf{Publicadas} \\
\textit{142/14} & \textit{Ríos Mojotoro y La Caldera, matrícula 3616} & \textit{Remitidas a la sede} \\
\textit{236/14} & \textit{Rectifica parcialmente la 142/14} & \textit{Remitidas a la sede} \\
182/14 & Río La Caldera, matrícula 2882 & \textbf{Publicadas} \\
\textbf{185/14} & \textbf{Río La Caldera y arroyo El Manzano, matrícula 4366} & \textbf{Publicadas} \\
209/14 & Río Vaqueros, matrículas 1905 y 3740 & \textbf{Publicadas} \\
\midrule
\textbf{383/14} & \textbf{Arroyos El Durazno y Guaranguay, \textit{matrícula 4366}} & \textbf{Remitidas a la sede} \\
\bottomrule
\end{longtable}
}

\begin{cautela}
\textbf{En ese tramo, nueve de once publicaron sus coordenadas} ---con una reserva sobre el numerador: las coordenadas que se publicaron para la Resolución 6/14 y para la 7/14 son las de la línea del pueblo que después aprobó la 023/14, según se explica más abajo, y al menos las de la 7/14 no corresponden a su inmueble. Contando sólo las que publicaron las propias, la proporción es ocho de once, o siete si tampoco las de la 6/14 son suyas, y el sentido no cambia---. Las dos que no lo hicieron son la Resolución 142/14, sobre los ríos Mojotoro y La Caldera en la matrícula 3616, y la Resolución 383/14, sobre los arroyos del loteo que se inundó cuatro años después.

La 142/14, además, tuvo que ser corregida: por \textbf{Resolución 236 del 25 de agosto de 2014} la Secretaría rectificó parcialmente su artículo 1º, y avisó que ``las coordenadas correctas se encuentran a disposición de los interesados en la sede''. Una determinación con coordenadas equivocadas, rectificada sin publicar ni las erróneas ni las correctas. Los dos avisos lo confirman: el de la 142/14 salió el 13 y el 14 de agosto de 2014 ---Boletines Nº 19.363 y 19.364, orden 100042974---, dos meses y medio después de dictada la resolución; el de la 236/14, el 3 de septiembre ---Boletín Nº 19.377, orden 100043378---, por un solo día. Ninguno ofrece enlace «Ver anexo». El segundo, además, cambia de régimen de impugnación: en lugar de la apelación ante el Tribunal de Aguas invoca el artículo 309 del Código de Aguas y el Decreto 1100, y el índice del Boletín le asigna el expediente «34-19.223/09» en lugar del 125-19.223/09 de la resolución que rectifica.

\textbf{Corresponde una rectificación.} Una versión anterior de este apartado afirmaba que la 383/14 era la única de la serie sin coordenadas publicadas. \textbf{No lo es}: la 142/14, de mayo de 2014, usa exactamente la misma fórmula --- ``se encuentran a disposición de los interesados en la sede'' --- y también remite. La afirmación de unicidad era errónea y queda retirada.

Lo que subsiste, y es lo que importa, es que \textbf{la práctica dominante fue publicar}: nueve de once en 2013--2014. Las tres determinaciones de 2015 sobre la Finca Wierna ---la 293/15 y la 428/15--- y sobre el río Vaqueros ---la 322/15--- también publicaron sus coordenadas; la cuarta de ese año, la 467/15, no, y se analiza enseguida. Con ellas la serie queda en doce de quince, y la proporción se sostiene.
\end{cautela}

\begin{cautela}
Y las dos excepciones tienen algo en común.

En octubre de 2015 la Secretaría de Recursos Hídricos publicó la gestión de \textbf{El Chalchal S.R.L.}, ``co-titular registral del inmueble \textbf{Catastro Nº 3616} del Dpto.\ La Caldera'', para obtener concesión de agua subterránea destinada al \textbf{abastecimiento poblacional de 580 habitantes --- 116 lotes --- del Loteo ``La Misión Etapa III''}. La matrícula 3616 es la de la Resolución 142/14; la 4366, la de la 383/14, sobre la que se tramitaba el agua para \textbf{cuatrocientos lotes de El Durazno}.

\textbf{Las dos determinaciones de ese tramo que no publicaron coordenadas recaen sobre suelo que este relevamiento encontró en proceso de loteo.} Pero la condición no basta para explicar la diferencia: entre las nueve que publicaron hay dos que también recaen sobre suelo en loteo ---la 80/14, sobre el inmueble del club de campo Las Vertientes, y la 185/14, sobre la misma matrícula 4366---, de modo que la coincidencia es parcial y no autoriza una regla.

\textbf{Hay un tercer caso, y acota la afirmación.} La Resolución 467/15, de diciembre de 2015, aprobó la línea de ribera del \textbf{río Lesser} entre dos puntos geográficos y también remitió las coordenadas a la sede. Es un acto de otra clase ---delimita un tramo de curso y no un inmueble--- y su pertenencia al departamento es probable ---los edictos de 2009 y 2011 que este capítulo transcribe ubican los ríos Lesser y Castellanos en el departamento La Caldera---, aunque el paraje Lesser aparece en otros actos bajo jurisdicción de San Lorenzo. Por eso la distinción se formula entre determinaciones \textbf{referidas a inmuebles}.

Y hay un dato que descarta la explicación benigna. \textbf{La Resolución 185/14 recae sobre la misma matrícula 4366 que la 383/14}: mismo inmueble ---la Finca La Caldera o Getsemaní---, mismo organismo, mismo asesor legal firmante, cinco meses de diferencia. Más aún, las dos salen del mismo expediente. En julio de 2014 ese edicto ---Boletines Oficiales Nº 19.346 y 19.347, del 21 y el 22 de julio--- publicó \textbf{ciento trece pares de coordenadas} ---trazados en la lámina \ref{fig:formacionradios}--- de los otros dos cursos de la finca: veintiuno del río La Caldera, catorce en la margen derecha y siete en la izquierda, y noventa y dos del arroyo El Manzano, en dos tramos. Lo hizo con un arancel de doscientos cincuenta pesos; en diciembre, sobre El Durazno y Guaranguay, remitió las suyas a la sede. El edicto de la Resolución 348/13 había publicado once pares ---u ocho, si se cuentan sólo los de la línea; la diferencia sigue sin cotejarse--- por ciento ochenta pesos: \textbf{entre diez y catorce veces más coordenadas por un cuarenta por ciento más de arancel}. La hipótesis del costo, de la composición o de la extensión no se sostiene.

Este libro \textbf{no afirma} que la omisión haya sido deliberada, ni que alguien la decidiera para favorecer a nadie: puede ser una decisión de trámite, un criterio del profesional interviniente o un error, y con once casos las coincidencias ocurren. Lo que hace es registrar el hecho y formular la pregunta: \textbf{por qué las coordenadas de los arroyos que atraviesan el loteo El Durazno no se publicaron, cuando las de los otros dos cursos de la misma finca sí.} Se responde con los dos expedientes y con el informe técnico de sus comisiones, y ninguno está en poder de este trabajo.
\end{cautela}

\subsection*{Ocho años después, el organismo dibujó el resultado}

Todo lo anterior se reconstruyó leyendo edictos de a uno. En mayo de 2022, cumpliendo un
mandato de la sentencia que analiza el capítulo \ref{cap:amparo}, la propia Secretaría de Recursos Hídricos
puso esa serie en un mapa: la \textbf{lámina \ref{lam:mapa2}} es la ubicación de las líneas de
ribera registradas por el organismo sobre el río La Caldera y sus cauces aportantes, aguas
arriba de la confluencia con el Wierna, a escala \textbf{1:40.000} y en \textbf{POSGAR 2007,
Faja 3}.

Vale la pena mirar qué trae, además de los trazos de ribera. La misma hoja superpone los
\textbf{polígonos catastrales} clasificados en urbano, subrural y rural, y los topónimos de las
fincas que este libro viene siguiendo desde el capítulo \ref{cap:fincas}: El Durazno, Guaranguay, Finca de
Linares, El Manzano, La Calderilla, Wierna, Los Yacones. \textbf{Las líneas de ribera y las
parcelas están, por primera vez, en el mismo plano y en el mismo sistema de coordenadas.}

\lamina{lam-mapa2-ribera.jpg}{lam:mapa2}%
{Líneas de ribera registradas sobre el río La Caldera y sus afluentes, hasta 2022}%
{\textbf{Mapa Nº 2.} Ubicación de las líneas de ribera en márgenes del río La Caldera y sus
cauces aportantes, aguas arriba de la unión con el río Wierna, registradas en la Secretaría de
Recursos Hídricos hasta el año 2022. Escala 1:40.000, sistema POSGAR 2007 --- Faja 3. Sobre la
misma base, los polígonos catastrales clasificados en urbano, subrural y rural.}%
{«Informe de mapas de cuenca de río La Caldera», Secretaría de Recursos Hídricos de Salta, mayo
de 2022, presentado en la ejecución de sentencia del amparo colectivo, expediente
INC-23860/1. Reproducción facsimilar; se recortó el margen en blanco del escaneo y no se
alteraron los tonos.}

\begin{cautela}
Dos advertencias sobre el uso de esta lámina.

\textbf{No reemplaza el cotejo.} A la resolución del ejemplar incorporado al expediente puede
verse dónde hay trazos de ribera y dónde no, pero no puede establecerse con seguridad qué
resolución corresponde a cada trazo ni si el inventario gráfico coincide exactamente con los
veinticinco actos que este capítulo reconstruyó. Eso se resuelve con las capas vectoriales, que
la Secretaría no publica; la lámina \ref{fig:riberaedictos}, más abajo, reconstruye una parte
con las coordenadas que sí publicó.

\textbf{Y no cubre el departamento entero.} El mapa se detiene en la confluencia con el río
Wierna: las determinaciones sobre el río Vaqueros --- la 209/14, la 322/15 --- y los arroyos
del municipio vecino quedan fuera del recorte. La partición cartográfica repite, sin
proponérselo, la partición jurisdiccional que el capítulo \ref{cap:vaqueros} documenta.
\end{cautela}

\subsection*{Las líneas que se pueden dibujar}

Veinte resoluciones de la serie publican sus coordenadas completas: diecinueve en el cuerpo del aviso y la 023/14 en el anexo vinculado a él. Las dos que publicaron coordenadas ajenas ---la 6/14 y la 7/14--- se explican más abajo y no se dibujan. Transcriptos y llevados al mismo sistema que la capa catastral, permiten dibujar lo que el Mapa Nº~2 muestra sin identificar: qué línea corresponde a cada resolución (lámina \ref{fig:riberaedictos}).

\begin{figure}[p]
\centering
\includegraphics[width=\textwidth,height=0.80\textheight,keepaspectratio]{fig-ribera-edictos.png}
\caption[Las líneas de ribera con coordenadas publicadas]{\textbf{Las líneas de ribera con coordenadas publicadas.} Los puntos de vinculación de veinte resoluciones de la Secretaría de Recursos Hídricos (dictadas entre 2012 y 2017; las cuatro del arroyo Chaile, en un solo color), unidos en el orden de su numeración, sobre la capa catastral. Línea continua: línea de ribera; punteada: línea de zona inundable. A: Vaqueros; B: el río Wierna y La Calderilla; C: el pueblo de La Caldera. No se dibujan las Resoluciones 6/14 y 7/14, cuyos avisos publicaron la tabla de la 023/14. Elaboración propia con las coordenadas de los avisos del Boletín Oficial ---ediciones y órdenes de publicación en la ficha y en la planilla del repositorio cartográfico---, tomadas como Gauss-Krüger faja 3 sobre POSGAR 94 (EPSG:22183) cuando el aviso no declara el marco y leídas según su magnitud cuando los rótulos X e Y están invertidos; la capa «Catastro Urbano y Rural» (IDESA, septiembre de 2026, orientativa) y la hidrografía del IGN. La diferencia entre POSGAR 94 y POSGAR 2007, marco del Mapa Nº~2, es inferior a un metro a esta escala. Elaboración propia; dibujo de \texttt{scripts/mapa\_ribera.py} del repositorio \texttt{mapas\_caldera} (apéndice \ref{cap:fuentes}).}\label{fig:riberaedictos}
\end{figure}

\begin{ficha}{Coordenadas publicadas, edicto por edicto}
\begin{itemize}[leftmargin=*,nosep]
\item \textbf{Res.\ 94/12}, río Castellanos, matrícula 4027: 4 vértices (V9, V10, V13, V14); el edicto rotula «Longitud» y «Latitud» columnas que contienen coordenadas planas.
\item \textbf{Res.\ 201/13}, río Vaqueros, sección B: 13 puntos en la margen izquierda y 5 en la derecha, «Coordenadas Posgar».
\item \textbf{Res.\ 348/13}, río La Caldera, matrícula 1590: 4 y 4 puntos, más 3 de zona inundable.
\item \textbf{Res.\ 80/14}, arroyo Quintín o Pacará, matrícula 3989: 20 y 16 puntos, «Gauss Kruger».
\item \textbf{Res.\ 185/14}, matrícula 4366: 7 y 14 puntos sobre el río La Caldera, 43 y 49 sobre el arroyo El Manzano. \textbf{Son 113, la cifra que da este capítulo.}
\item \textbf{Res.\ 428/15}, río La Caldera, matrícula 2020: 13 y 4 puntos.
\item \textbf{Res.\ 14/16}, río Wierna, matrícula 4445: 3 y 3 puntos, más 4 de línea de inundación; columnas rotuladas «Y, X», con la coordenada este bajo «Y».
\item \textbf{Res.\ 144/17}, río Vaqueros, matrícula 40 (ampliación): 4 y 4 puntos; el edicto llama X a la coordenada norte e Y a la este.
\item \textbf{Res.\ 023/14}, río La Caldera a lo largo del pueblo, sin matrícula: 58 puntos en la margen derecha y 76 en la izquierda, en el anexo vinculado a los avisos de los B.O.\ Nº 19.250 y 19.251 (orden 100039394); columnas «Y, X» con la norte bajo «Y».
\item \textbf{Res.\ 101/14}, arroyo Los Nogales Norte o El Cajón, matrícula 4191 (B.O.\ Nº 19.299, orden 100041196): 23 y 37 puntos en el cauce principal y 17 y 13 en un brazo que el aviso titula «Los Nogales Sur»; un bloque dice «Matrícula N° 491», el punto LRI-8 se repite y el LRI-50 y el LRD-50 tienen las mismas coordenadas. Llama X a la norte.
\item \textbf{Res.\ 182/14}, río La Caldera, matrícula 2882, La Calderilla (B.O.\ Nº 19.347, orden 100042593): 3 y 3 puntos, con rótulos opuestos en cada margen.
\item \textbf{Res.\ 209/14}, río Vaqueros, matrículas 1905 y 3740 (B.O.\ Nº 19.356, orden 400005047): 8 y 6 puntos, «Gauss Kruger».
\item \textbf{Res.\ 293/15}, matrículas 3721 y 3722, Finca Wierna (B.O.\ Nº 19.666, orden 400007168): 10 y 11 puntos sobre el río La Caldera, 8 y 7 sobre la Cañada Oeste. Llama X a la norte.
\item \textbf{Res.\ 322/15}, río Vaqueros, matrículas 156, 23, 1695 y 406 (B.O.\ Nº 19.675, orden 100051198): 2 y 4 puntos; columnas «Y, X» con la este bajo «Y».
\item \textbf{Res.\ 414/16}, río Vaqueros, matrícula 40 (B.O.\ Nº 19.931, orden 100057811): 4 y 3 puntos. Llama X a la norte.
\item \textbf{Res.\ 154/16}, río Wierna, matrícula 3722 (B.O.\ Nº 20.136, orden 400010474): 5 y 5 puntos, redondeados al metro. Llama X a la norte.
\item \textbf{Res.\ 205, 210, 211 y 212/17}, arroyo Chaile, parcelas 32 a 29 de la manzana 148 de Vaqueros (B.O.\ Nº 20.249 y 20.253): dos puntos por margen cada una, cuatro decimales, columnas «Y, X» con la este bajo «Y». Un punto de la 205/17 (X = 726931,0916) no se dibuja.
\end{itemize}
\textit{Quinientos treinta y cinco puntos en total, contando una vez por resolución los puntos de borde que dos resoluciones comparten. Transcripción del texto de cada edición, cotejada contra la imagen de la página donde la capa de texto era dudosa. La planilla, con edición, página, orden de publicación y fecha de consulta por punto, integra el repositorio cartográfico del libro (\texttt{datos/ribera/}).}
\end{ficha}

\begin{cautela}
\textbf{Lo que el dibujo confirma, lo que muestra y lo que no puede mostrar.}

\textbf{Confirma dos cifras del libro.} La Resolución 185/14 publica 113 puntos, y el último punto de la margen derecha de la 201/13 queda a trece metros y medio del último de la ampliación de la 144/17: la continuidad del río Vaqueros que este capítulo afirmaba se ve en el plano.

\textbf{Muestra que el marco de referencia no es uniforme.} De las veinte resoluciones, diez no declaran el sistema y una rotula como geográficas coordenadas planas. Once llaman X a la coordenada norte ---la convención geodésica---, ocho a la este, y la 182/14 usa una convención distinta en cada margen del mismo aviso. La 154/16 publica valores redondeados al metro; las del Chaile, con cuatro decimales. Con las coordenadas publicadas, cualquiera puede reproducir el trazo; con los rótulos tal como están, no sin interpretar.

\textbf{Y no muestra lo que falta.} No están aquí la 383/14 ni la 142/14 y su rectificación 236/14, cuyas coordenadas se remitieron a la sede, ni las determinaciones que, desde la 467/15 de diciembre de 2015, hicieron lo mismo. Tampoco la 6/14 ni la 7/14, porque lo que sus avisos publicaron es la línea del pueblo, que está dibujada una sola vez, como 023/14. Lo que se ve es todo lo que la serie publicó con coordenadas, no todas las líneas que el organismo fijó.
\end{cautela}

\begin{cautela}
Un apunte toponímico, registrado sin conclusión. El arroyo cuya línea de ribera se determinó en abril de 2014 se denomina \textbf{Los Nogales Norte o El Cajón}, y atraviesa la Finca Vaqueros. Los loteos afectados por los aluviones de 2018 y 2019 en el municipio vecino se llaman \textbf{El Nogalar} I, II y IV.

La proximidad de los nombres puede indicar continuidad de un mismo accidente geográfico o simple repetición de un topónimo frecuente en la zona --- los nogales abundan en las yungas del corredor. Establecerlo requiere cartografía, no lectura. Se consigna porque, si se tratara del mismo sistema hídrico, la determinación de abril de 2014 alcanzaría a un curso vinculado a los loteos que se inundaron.
\end{cautela}

\begin{cautela}
Dos advertencias documentales, ninguna resuelta.

\textbf{Sobre el expediente.} El sumario del Boletín Oficial Nº 19.466 consigna, para el aviso de la Resolución 383/14, el expediente \textbf{34-167.127/13}, el mismo de la secuencia 395/13--185/14. \textbf{Un solo expediente administrativo cubre la finca entera y sus cuatro cursos de agua}: de él salieron la 395/13, que designó la comisión técnica; la 185/14, que determinó la línea del río La Caldera y del arroyo El Manzano; y la 383/14, la de El Durazno y Guaranguay. Tres actos, un expediente, una matrícula ---la 4366, pese a que el aviso de la 395/13 dice ``1366'' (B.O.\ Nº 19.232, 24/01/2014): la errata está en el aviso mismo, no en la lectura del escaneo.

\textbf{Sobre las coordenadas.} Los avisos de las Resoluciones 6/14 y 7/14 (B.O.\ Nº 19.232 y 19.233, 24 y 27/01/2014; órdenes 100038934 y 100038932) publican la misma tabla de ciento treinta y cuatro puntos, y \textbf{esa tabla es, punto por punto, la del anexo de la Resolución 023/14}: la línea del río La Caldera a lo largo del pueblo, cincuenta y ocho puntos en la margen derecha y setenta y seis en la izquierda, cotejados uno a uno contra el facsímil del anexo. Tres cosas se siguen.

\textbf{La tabla no se arrastró de la 023/14.} Los dos avisos salieron entre dos y tres semanas antes de que la 023/14 se dictara (06/02/2014) y casi un mes antes de que se publicara (19/02/2014). La tabla existía antes: es, con toda probabilidad, la determinación que la Dirección General de Planificación Hídrica hizo «por Administración» y que la 023/14 cita a fojas 3 de su expediente.

\textbf{Al menos la 7/14 publicó coordenadas ajenas.} Los ciento treinta y cuatro puntos van, de norte a sur, de 7.281.221 a 7.277.907 metros: el tramo del pueblo. La 7/14 recae sobre la matrícula 3584 de La Calderilla, y la determinación publicada más próxima en esa localidad ---la 182/14, matrícula 2882--- está unos dos kilómetros y medio más al sur que el último punto de la tabla. Dónde está la matrícula 3593 de la 6/14 este libro no lo sabe; si no está en el tramo del pueblo, sus coordenadas tampoco son las suyas. La capa catastral lo resolvería si la matrícula sigue vigente.

\textbf{Y la línea del pueblo quedó publicada tres veces bajo tres actos distintos}: dos que se refieren a inmuebles particulares y fueron determinados por comisiones técnicas ---la de la 6/14 la designó la Resolución 322/12---, y uno pedido por el intendente y hecho por la propia Administración. Qué línea rige para las matrículas 3593 y 3584 es una pregunta que sólo el organismo puede contestar; el apéndice \ref{cap:pedidos} la formula.
\end{cautela}

\begin{cautela}
Un quinto acto del mismo período cierra el círculo de otra manera.

Los Boletines Nº 19.250 y 19.251, del 19 y el 21 de febrero de 2014, publican la \textbf{Resolución Nº 000023/14} de la Secretaría de Recursos Hídricos con un título ---«Río La Caldera, Pueblo La Caldera»--- y la leyenda «ver anexo». El anexo, vinculado a los dos avisos en la edición web, contiene la resolución completa.

\end{cautela}

\begin{ficha}{Resolución SRH Nº 000023 del 6 de febrero de 2014 --- expte.\ 0090034-96837/2013}
\textbf{Origen:} nota del intendente de La Caldera, \textbf{Luis G.\ Mendaña\index{Mendaña, Luis Gerardo}}, que pide la línea de ribera de ambas márgenes del río La Caldera a lo largo del pueblo, con una imagen satelital (fs.\ 1 y 2). La Dirección General de Planificación Hídrica informa que la determinación \textbf{«se ha realizado por Administración»}, con coordenadas y plano, y que \textbf{«estará sujeta a la construcción de defensas con gaviones a lo largo de la misma»}, tramitada por expediente 34-253.934/13 (fs.\ 3). Dictamen jurídico por el art.\ 126 del Código de Aguas y el Decreto 1989/02 (fs.\ 7).

\textbf{Art.\ 1:} aprueba la designación del \textbf{Geól.\ Jorge G.\ Torres\index{Torres, Jorge G.}} para las tareas de determinación.

\textbf{Art.\ 2:} \textbf{determina la línea de ribera sobre ambas márgenes del río La Caldera a lo largo del pueblo}, con sus coordenadas: \textbf{134 puntos}, 58 en la margen derecha y 76 en la izquierda.

\textbf{Art.\ 3:} \textbf{advierte a la Municipalidad que la determinación estará sujeta a la construcción de defensas con gaviones} a lo largo de ella.

Arts.\ 4 y 5: publicación por dos días y notificación a Inmuebles. Firma el secretario Alfredo Fuertes.
\end{ficha}

\begin{cautela}

\textbf{La primera:} Jorge G.\ Torres\index{Torres, Jorge G.} es el autor de la presentación institucional sobre línea de ribera que este capítulo utiliza como fuente del procedimiento, y figura en diciembre de 2025 como Director General de Cuencas y Política Hídrica de la provincia. En febrero de 2014 fue designado para determinar la línea de ribera del río La Caldera \textbf{a la altura del pueblo}. La continuidad entre quien ejecutó la tarea en el terreno y quien hoy conduce el área es completa.

La lámina \ref{fig:otbn} traza esa línea sobre el pueblo.

\textbf{La segunda, y es una corrección.} Este libro había registrado esta resolución como publicada sin contenido y la contaba entre los anexos no publicados. No lo es: el anexo está vinculado a los dos avisos y trae las coordenadas completas. \textbf{El acto más pertinente al ejido urbano de La Caldera es, al contrario, uno de los que publicaron sus coordenadas}, aunque en un formato ---una copia escaneada vinculada a un aviso de dos líneas--- que ninguna búsqueda de texto encuentra.

\textbf{La tercera: la línea de ribera del pueblo quedó condicionada a una obra.} El artículo 3 sujeta la determinación a defensas con gaviones que en febrero de 2014 se tramitaban por un expediente propio. La línea se trazó contando con defensas que todavía no existían, y es ese expediente ---no la resolución--- el que dice si se construyeron (capítulo \ref{cap:defensas}). Y el pedido lo hizo el intendente: el municipio promovió la determinación sobre su propio pueblo.
\end{cautela}

Y la matrícula 4366 gana un nombre y una dimensión. Se llama \textbf{Finca La Caldera o Getsemaní}, y está atravesada por al menos \textbf{cuatro cursos de agua}: el río La Caldera, el arroyo El Manzano, el arroyo El Durazno y el arroyo Guaranguay. Las dos determinaciones de 2014 --- julio y diciembre --- cubren dos cursos cada una.

Eso reordena la pregunta dominial que este libro arrastra. La Resolución 646D/21 aprobó el plano del loteo El Durazno sobre la \textbf{matrícula 6.558}. Las líneas de ribera de sus arroyos se determinaron sobre la \textbf{matrícula 4366}. Si la 6.558 proviene de la 4366 por desmembramiento --- lo que la numeración sugiere y nada prueba ---, entonces el loteo se asienta sobre una finca cuyos cuatro cursos de agua fueron relevados por el Estado entre 2013 y 2014, y la línea de ribera de dos de ellos recae directamente sobre su suelo.

Es la comprobación más importante que le queda a este capítulo, y se hace con un informe de dominio de la Dirección General de Inmuebles.

\begin{cautela}
\textbf{Después de la serie, el mismo patrón.} La serie de 2013 a 2015 no fue la última palabra. El Boletín Oficial Nº 19.951, del 27 de enero de 2017, publica el edicto de la \textbf{Resolución Nº 360 del 7 de noviembre de 2016} de la Secretaría de Recursos Hídricos, expediente 34-171.172/15.

\end{cautela}

\begin{ficha}{Resolución SRH Nº 360/16 --- B.O.\ Nº 19.951, pág.\ 71}
Aprueba la determinación de la línea de ribera realizada por la Comisión Técnica en \textbf{ambas márgenes del río Wierna}, en la zona de los inmuebles de \textbf{matrículas 4097, 4953, 391, 389, 1268, 3718 y 3438} del departamento La Caldera. Integran la comisión el geólogo \textbf{Jorge Genaro Torres} y el ingeniero civil Carlos Enrique Cerezo.

``Las coordenadas Posgar de las líneas de ribera determinadas en ambas márgenes de los cauces mencionados \textbf{se encuentran a disposición de los interesados en la sede} de la Secretaría de Recursos Hídricos.''

Apelable ante el Tribunal de Aguas en diez días hábiles. Firman el secretario Alfredo Fuertes y la asesora legal Sandra Mabel Siegrist; edicto fechado el 16 de noviembre de 2016 y publicado el 26 y el 27 de enero de 2017.
\end{ficha}

\begin{cautela}

Tres observaciones. La fórmula es literalmente la de las Resoluciones 142/14 y 383/14, con otro firmante, dos años después: la omisión de las coordenadas no fue un rasgo de un asesor ni de una racha. El geólogo que la Resolución 023/14 designó para el río La Caldera sigue integrando comisiones del departamento en 2016. Y el río Wierna ---donde el capítulo \ref{cap:amparo} registra cuatro canteras y el capítulo \ref{cap:defensas} un encauzamiento provincial en 1980--- tiene línea de ribera determinada desde 2016 sobre siete inmuebles cuya relación con esas canteras y obras este libro no estableció.
\end{cautela}

\section{2017: la línea de ribera que pide el propio Estado}

Este capítulo miró hasta acá quién pide una determinación de línea de ribera: un
peticionante particular, un loteo, un vecino. Falta un tercer solicitante, y es el más
revelador.

\begin{ficha}{Secretaría de Comunicación de la Provincia de Salta, 16 de agosto de 2017}
El secretario de Recursos Hídricos, \textbf{Alfredo Fuertes}, se reúne con profesionales de la
empresa \textbf{Supercemento}, la consultora \textbf{IHTSA} y la empresa \textbf{Aguas del
Norte} para coordinar las tareas de la obra del \textbf{Acueducto Norte}.

Los participantes destacan como punto principal \textbf{la determinación de la línea de ribera
del río La Caldera en la zona de emplazamiento de la planta de tratamiento}, porque de ese
trabajo depende que puedan diseñarse las defensas apropiadas. Se explican además los
requisitos necesarios para los cruces de los cauces de \textbf{Los Naranjos, Guaranguay y
Wierna}.

El acueducto tendrá una capacidad de \textbf{1.100 metros cúbicos por hora} y abastecerá a
toda la zona norte de la ciudad de Salta y a los pobladores de Vaqueros.

Participa, por la Provincia, el \textbf{director de Planificación Hídrica, Jorge Torres}.
\end{ficha}

\textbf{Primera: aparece Guaranguay otra vez.} Es el mismo arroyo cuya línea de ribera se
determinó en diciembre de 2014 sin publicar coordenadas. Dos años y ocho meses después, el
propio Estado necesita saber dónde está esa línea para cruzarlo con un acueducto. El dato que
no se publicó para el vecino hacía falta para la obra.

\textbf{Segunda: el instrumento sirve, y el Estado lo sabe.} Cuando la obra es del Estado, la
determinación de línea de ribera no es un trámite ni un requisito formal: es el paso previo
sin el cual no pueden diseñarse las defensas. Lo dice la propia comunicación oficial. Es la
mejor prueba disponible de que el instrumento no es decorativo, y de que su ausencia en El
Durazno no se explica por irrelevancia técnica.

\textbf{Y tercera: Jorge Torres cambia de lugar en la estructura.} En febrero de 2014 la
Resolución 023/14 lo designa \textbf{como geólogo} para determinar la línea de ribera del río
La Caldera a la altura del pueblo. En agosto de 2017 participa \textbf{como director de
Planificación Hídrica}. Se consigna el recorrido porque explica por qué su nombre aparece en
los dos extremos de este capítulo, y no como imputación de ninguna clase.

\begin{cautela}
La fuente es \textbf{un parte de prensa oficial}, no un expediente. Acredita que la reunión
existió y qué se dijo en ella; \textbf{no acredita que la determinación se haya hecho}. Este libro no encontró la resolución de determinación, y la duda queda abierta: si el Estado determinó esa línea, el acto debería existir; si no la determinó, la obra se diseñó sin ella. Lo que sí encontró es una comisión técnica de 2015 sobre el río La Caldera en el catastro 4049, con Torres entre sus integrantes, cuyo aviso el índice del Boletín atribuye a Aguas del Norte (véase más adelante en este capítulo).
\end{cautela}

\section{Trece meses antes, el mismo organismo hizo lo contrario}

Existe un edicto anterior del mismo organismo, sobre el mismo río, en el mismo departamento y con la misma firma. Sirve de control y convierte una impresión en una comparación.

\begin{ficha}{Resolución Nº 348 del 19/11/2013 --- B.O.\ Nº 19.213, 23 y 26 de diciembre de 2013}
Determinación de la línea de ribera en ambas márgenes del \textbf{río La Caldera}, en la zona del inmueble de \textbf{matrícula catastral Nº 1590, localidad de La Calderilla}, departamento La Caldera.

\textbf{El edicto publica las coordenadas.} Cuatro puntos Posgar de margen izquierda (LRi01 a LRi04) y cuatro de margen derecha (LRd05 a LRd08), con sus valores X e Y completos. \textit{(Más adelante este capítulo cuenta once pares publicados en este edicto; la diferencia con los ocho de la línea de ribera debe cotejarse contra el original.)}

\textbf{El edicto determina además la zona inundable}, con su propia tabla de coordenadas, y advierte que ``puede sufrir inundaciones durante las crecidas extraordinarias''.

\textbf{Y transcribe el artículo 172 del Código de Aguas}, con sus restricciones al uso del suelo.

Firma: Dr.\ Matías J.\ Brogin\index{Brogin, Matías J.}, Asesor Legal. Orden de publicación 400004557.
\end{ficha}

\begin{cautela}
Póngase al lado el edicto de la Resolución 383/14, publicado trece meses después, en enero de 2015, por la misma Secretaría y con la firma del mismo asesor legal, sobre los arroyos El Durazno y Guaranguay.

\begin{center}\small
\begin{tabular}{@{}p{4.4cm}p{3.8cm}p{3.8cm}@{}}
\toprule
& \textbf{Res.\ 348/13} (dic.\ 2013) & \textbf{Res.\ 383/14} (ene.\ 2015) \\
\midrule
Coordenadas Posgar & \textbf{Publicadas en el edicto} & Remitidas a la sede \\
Zona inundable & \textbf{Determinada y publicada} & No mencionada \\
Art.\ 172 del Código de Aguas & \textbf{Transcripto} & No mencionado \\
\bottomrule
\end{tabular}
\end{center}

Este libro había observado que el edicto de 2015 remitía las coordenadas a Avenida Bolivia 4650 y que eso volvía difícil el ejercicio del derecho a apelar. La observación se sostenía sobre el sentido común. Ahora se sostiene sobre un precedente: \textbf{el organismo sí publicaba coordenadas, y lo había hecho poco antes, en el mismo departamento}. No era una imposibilidad técnica ni una práctica uniforme.

\textbf{Corresponde una advertencia de peso.} El edicto es una síntesis de la resolución, no la resolución. Que el edicto de 2015 no mencione la zona inundable ni el artículo 172 \textbf{no prueba} que la Resolución 383/14 no los haya determinado o invocado. Puede haberlo hecho y el edicto haber salido más escueto. Establecerlo requiere el texto de la resolución, no su publicación.

Lo que queda probado es más acotado y sigue siendo sustancial: \textbf{lo que llegó a conocimiento público fue distinto en un caso y en otro}, y el caso que llegó mermado es el que recae sobre los arroyos que se desbordaron cuatro años después.
\end{cautela}

\subsection*{Y en 2016, con la misma firma, volvió a publicarlo todo}

\begin{ficha}{Resolución Nº 14 del 15/01/2016 --- B.O.\ Nº 19.740 y 19.741, 15 y 16 de marzo de 2016}
Expediente 34-125.733/15. Aprueba la determinación de la línea de ribera realizada por la Comisión Técnica en \textbf{ambas márgenes del río Wierna}, en la zona de la matrícula 4445, Fracción 195, Finca «Vaqueros», localidad Vaqueros.

\textbf{Publica las coordenadas}: tres puntos Gauss-Krüger por margen (LRI-01 a 03 y LRD-01 a 03).

\textbf{Determina una línea de inundación} sobre la margen derecha, con cuatro puntos (Li-01 a Li-04), declara que esas zonas «pueden sufrir inundaciones durante las crecidas extraordinarias» y \textbf{transcribe el artículo 172 del Código de Aguas}.

Edicto del 25 de enero de 2016. Firman el secretario Alfredo Fuertes y el asesor legal \textbf{Matías J.\ Brogin}\index{Brogin, Matías J.}.

La comisión se había aprobado por la Resolución 273/15 del 13 de agosto de 2015 (B.O.\ Nº 19.614 y 19.615) para «ambas márgenes del \textbf{río Trancas}», con el ingeniero Antonio Alejandro Forns\index{Forns, Antonio Alejandro}, el geólogo Teodoro Chafatinos y el ingeniero Lucas Vidal.
\end{ficha}

\begin{cautela}
Este edicto cambia el peso de la comparación anterior. La Resolución 348/13 no fue un caso único: un año después de publicar sin coordenadas ni zona inundable la línea de El Durazno, el mismo organismo, con el mismo asesor legal, publicó para un inmueble de Vaqueros las coordenadas, la línea de inundación y el artículo 172. La forma completa del edicto existía antes y después de la Resolución 383/14. Lo que la diferencia no es el período ni el firmante.

Y deja una discordancia que este libro registra sin resolver: la comisión se designó para el río Trancas y la determinación se aprobó sobre el río Wierna, en el mismo expediente y el mismo inmueble. Puede tratarse de dos nombres del mismo tramo o de un error de uno de los dos edictos. \pendiente{aclaración de la Secretaría sobre el curso del expediente 34-125.733/15}
\end{cautela}

\begin{ficha}{Resolución Nº 144/17 --- B.O.\ Nº 20.042 y 20.043, 15 y 16 de junio de 2017}
Expediente 34-202.169/15. Modifica parcialmente el artículo 1º de la \textbf{Resolución 414 del 5 de noviembre de 2016} ---fecha que la numeración de la serie contradice; véase la tabla--- y aprueba la \textbf{ampliación} de la determinación de la línea de ribera en ambas márgenes del \textbf{río Vaqueros}, en la zona que colinda con la \textbf{matrícula 40}. Publica cuatro puntos por margen (LRmd1a a 4a, LRmi5a a 8a). Edicto del 11 de mayo de 2017; firman Fuertes y la asesora legal Sandra Mabel Siegrist.
\end{ficha}

El expediente del río Vaqueros quedó completo con la comisión de la Resolución 147/16 y la determinación de la 414/16. Entre la margen derecha y la izquierda median unos ciento treinta a ciento cuarenta metros; con la ampliación de 2017, la margen derecha suma unos doscientos cuarenta metros de poligonal publicada. Y la comisión tiene una composición que ninguna otra de las leídas repite: tres miembros, uno de ellos \textbf{ingeniero hidráulico}, que es la especialidad que el artículo 3º del Anexo I pone en primer lugar.

Dos datos más para la serie. La omisión de las coordenadas tampoco fue un rasgo de la asesora que firma la serie desde 2016 ---y edictos del organismo al menos desde agosto de 2014, capítulo \ref{cap:redes}---: en noviembre de 2016 su edicto sobre el Wierna (Resolución 360/16) las remitió a la sede, y un mes después, en diciembre, su edicto sobre el río Vaqueros (Resolución 414/16) las publicó, como volvió a hacerlo seis meses más tarde con la ampliación (Resolución 144/17). Y una determinación publicada puede ampliarse después: la de la matrícula 40 se modificó a los seis meses. Qué motivó la ampliación, y si el tramo agregado coincide con una subdivisión o un proyecto sobre ese inmueble, es lo que diría el expediente. \pendiente{fundamentos de la ampliación de la Resolución 414/16 por la 144/17}

\subsection*{2016 y 2017: el relevamiento siguió}

La serie de 2013 a 2015 no se cerró en diciembre de 2015. Los avisos de 2016 muestran que el departamento siguió bajo relevamiento, con los mismos técnicos y dos cursos nuevos.

{\small
\setlength{\LTpre}{0.6em}\setlength{\LTpost}{1.2em}
\begin{longtable}{@{}p{2.3cm}p{2.6cm}p{7cm}@{}}
\toprule
\textbf{Acto} & \textbf{Expediente} & \textbf{Objeto} \\
\midrule
\endhead
Res.\ 154/16 (17/05/16) & 34-259.760/14 & \textbf{Determinación}: \textbf{río Wierna}, matrícula 3722, Finca Wierna ---el expediente de la comisión 74/15 y de la determinación 293/15 sobre el río La Caldera---. Con coordenadas, cinco puntos por margen \textbf{redondeadas al metro}. Edicto del 23/05/2016, \textbf{publicado el 2 y 3/11/2017}, diecisiete meses después \\
Res.\ 14/16 (15/01/16) & 34-125.733/15 & \textbf{Determinación}: río Wierna (comisión por el Trancas), matrícula 4445, Finca Vaqueros. \textit{Con coordenadas y línea de inundación} \\
Res.\ 122/16 (21/04/16) & 34-62.718/16 & \textbf{Comisión}: arroyo Quintín, matrícula 4953, Finca Vaqueros. Torres y Cerezo. B.O.\ 6 y 7/06/2016 \\
Res.\ 282/16 (26/08/16) & 34-202.924/15 & \textbf{Determinación}: río La Caldera, matrícula 1499, Finca San Jorge, La Calderilla. \textit{Coordenadas remitidas a la sede}; cita el «Dcto.\ 1898/02». B.O.\ 28 y 29/09/2016 \\
Res.\ 290/16 (07/09/16) & 34-178.741/16 & \textbf{Comisión}: \textbf{arroyo Chaile}, matrícula 1534. El ingeniero Marcelo Toigo y el geólogo Enrique Chalabe. B.O.\ 20 y 21/10/2016 \\
Res.\ 291/16 (07/09/16) & 34-116.088/16 & \textbf{Comisión}: río La Caldera, \textbf{matrícula 4146}, Fracción 71, La Calderilla. Torres y el ingeniero Emanuel Caro. B.O.\ 30/09 y 03/10/2016 \\
Res.\ 147/16 (11/05/16) & 34-202.169/15 & \textbf{Comisión}: río Vaqueros, catastro 40. Integran el ing.\ en construcciones Rolando Chávez, el \textbf{ingeniero hidráulico} Víctor Zacarías Pérez y el geólogo Héctor Saravia Navamuel. B.O.\ 28 y 29/06/2016 \\
Res.\ 414/16 (05/11/16) & 34-202.169/15 & \textbf{Determinación}: río Vaqueros, matrícula 40. \textit{Con coordenadas publicadas} (tres puntos en la margen derecha, cuatro en la izquierda). Fecha según la Res.\ 144/17: el edicto no la consigna. \textbf{La fecha no cuadra con la numeración}: la 360/16 es del 7 de noviembre y la 389/16 del 21, de modo que la 414 difícilmente sea del 5 de noviembre, que además fue sábado; el 5 de diciembre, lunes, encaja con la serie, pero este libro no lo establece. B.O.\ 29/12/2016 y 02/01/2017; ampliada por la 144/17 \\
Res.\ 389/16 (21/11/16) & 34-248.094/16 & \textbf{Determinación «realizada por Administración»}: arroyo Quintín, matrícula 4953, Finca Vaqueros. \textit{Coordenadas remitidas a la sede}; cita el «Dcto.\ 1898/02». Firma el geólogo Torres como director general de Planificación Hídrica. B.O.\ 19 y 20/12/2016 \\
Res.\ 360/16 (07/11/16) & 34-171.172/15 & \textbf{Determinación}: río Wierna, siete matrículas. \textit{Coordenadas remitidas a la sede} \\
Res.\ 424, 425, 426 y 427/16 (14/12/16) & 34-257.289/16, 34-229.971/16, 34-257.277/16 y 34-229.967/16 & \textbf{Cuatro comisiones}: arroyo Chaile o Teodolinda, matrículas 2343, 2345, 2342 y 2344 (Vaqueros, sección A, manzana 148, parcelas 30, 32, 29 y 31). Torres y Caro. B.O.\ Nº 19.934 y 19.935, 4 y 5/01/2017; dos de los títulos del índice dicen «río La Caldera», y el aviso de la 427/16 lo dice también en su cuerpo \\
Res.\ 205, 210, 211 y 212/17 & 34-229.971/16, 34-257.289/16, 34-229.967/16 y 34-257.277/16 & \textbf{Cuatro determinaciones}: arroyo Chaile o Teodolinda, matrículas 2345, 2343, 2344 y 2342. \textit{Con coordenadas publicadas}, dos puntos por margen cada una. Edictos de abril de 2018, publicados el 25 y 26/04 y el 3 y 4/05/2018 \\
Res.\ 279/17 (09/08/17) & 34-116.088/16 & \textbf{Determinación}: río La Caldera, matrícula 4146, Fracción 71, La Calderilla; cierra la comisión de la 291/16. \textit{Coordenadas remitidas a la sede}; cita el «Dcto.\ 1.898/02». Edicto del 18/08/2017, B.O.\ 20 y 21/09/2017 \\
\bottomrule
\end{longtable}
}

La Resolución 154/16 agrega un caso de demora extrema: firmada por Brogin en mayo de 2016, se publicó en noviembre de 2017, cuando los edictos del organismo ya los firmaba otra asesora. Y sus coordenadas, a diferencia de todas las demás publicadas, no tienen decimales: puntos separados por cuatrocientos cincuenta a setecientos cincuenta metros, redondeados al metro, a lo largo de más de dos kilómetros de río.

Todas las comisiones de 2016 invocan el método «geológico-geomorfológico»; las de 2015, «geológico-morfológico». Firman Brogin hasta mediados de año y Siegrist desde septiembre.

\begin{cautela}
Las actuaciones sobre el arroyo Chaile o Teodolinda en Vaqueros muestran algo que ningún otro grupo de la serie mostraba: \textbf{una sola línea de ribera partida en lotes}. Son cuatro parcelas contiguas de una manzana urbana, la 148 de la sección A, con matrículas correlativas (2342 a 2345) y parcelas correlativas (29 a 32). Cada una tiene su expediente, su comisión ---las cuatro del mismo día, con los mismos dos técnicos---, su determinación y sus dos edictos. Pero los puntos forman una sola poligonal: los de borde se repiten de una resolución a la siguiente con las mismas coordenadas. Cada tramo mide entre catorce metros y medio y dieciséis metros y medio por margen ---salvo la margen derecha de la parcela 32, que no se puede medir por el dígito faltante que se señala abajo---; los cuatro juntos, unos sesenta y dos.

La publicación tiene tres defectos menores que, juntos, muestran cómo se armaba. Los nombres de los puntos no son únicos: la parcela 31 y la 32 usan los mismos rótulos (LRmi01, LRmi02, LRmd09, LRmd10) para coordenadas distintas. Un punto compartido figura con un centímetro de diferencia entre dos edictos. Y una coordenada de la parcela 32 se publicó con un dígito de menos (X = 726931,0916, que no corresponde a ningún lugar de la provincia). A eso se suman dos títulos del índice que atribuyen el arroyo al río La Caldera.

Es, por lo tanto, un relevamiento continuo tramitado como cuatro. Qué lo explica ---cuatro propietarios distintos, un fraccionamiento que cada comprador debió completar por su cuenta, o la exigencia de una línea por lote para aprobar construcciones--- es lo que dirían los expedientes. Lo que el Boletín muestra es el resultado: ocho publicaciones pagadas para fijar poco más de sesenta metros de ribera. Las determinaciones llevan número de 2017 y se publicaron en abril y mayo de 2018. La comisión de la parcela 31 sí se publicó: es la Resolución 427/16, del mismo día que las otras tres y con su mismo expediente (B.O.\ Nº 19.934 y 19.935, 4 y 5/01/2017), pero su aviso la dicta para el «río La Caldera» y no para el arroyo Chaile. Una versión anterior de este párrafo la daba por no publicada. \pendiente{solicitantes de los cuatro expedientes}
\end{cautela}

\begin{cautela}
Tres cosas se siguen de esta tabla. La primera cierra una pregunta abierta: la serie del río La Caldera sobre la matrícula 4146 ---determinada por la Resolución 279/17, retomada por la 306/19 y rectificada por la 108/2025--- empieza con la Resolución 291/16, que designó la comisión en septiembre de 2016, en La Calderilla y con el geólogo Torres, y la determinación del 9 de agosto de 2017 la cierra en el mismo expediente, con las coordenadas remitidas a la sede, igual que la rectificación de 2025.  El edicto de la Resolución 306/19 no se leyó: la búsqueda por el número del expediente de 2016 devuelve sólo la comisión y la determinación de 2017, de modo que la 306/19 tramitó, probablemente, con un expediente de numeración nueva. La 279/17, en cambio, es anterior a la reunión de 2017 sobre el Acueducto Norte, y es la candidata más firme para la línea que esa reunión presentaba como pendiente.

La segunda es de nombres: el ingeniero Toigo, que en 2011 integró la comisión de Chacras del Gallinato, integra en 2016 la del arroyo Chaile. Y la matrícula 4953 aparece dos veces en el mismo año, en la comisión del Wierna y en la del Quintín: un mismo inmueble atravesado por dos cursos, con dos expedientes.

La determinación del Quintín de noviembre de 2016 agrega dos rasgos que la serie no tenía. No la hizo una comisión técnica sino, dice el edicto, «la Administración». Y la firma, junto a la asesora legal, el geólogo Jorge Genaro Torres como director general de Planificación Hídrica: el mismo técnico que integró la comisión designada en abril de ese año sobre el mismo arroyo y la misma matrícula, y que integra la mayoría de las comisiones del departamento entre 2015 y 2017, aparece aquí del lado de la autoridad que aprueba y publica. El expediente es otro (34-248.094/16, no el 34-62.718/16 de la comisión), de modo que no puede afirmarse que esta determinación cierre aquella comisión; puede ser un trámite de oficio distinto sobre el mismo tramo.

Hay además una anomalía de expediente. El edicto de la Resolución 131/16, que designa una comisión para el río San Lucas en la matrícula 559 del \textbf{departamento San Carlos}, lleva como referencia el expediente 34-62.718/16: el mismo de la comisión del arroyo Quintín, con la misma comisión (Torres y Cerezo), la misma fecha de edicto (29 de abril de 2016) y la misma firma. Uno de los dos encabezados está mal; es el mismo tipo de arrastre de texto entre edictos que la cita del decreto.

La tercera es de forma. En septiembre de 2016 la misma asesora firmó una determinación con las coordenadas remitidas a la sede y con el decreto mal citado; ocho meses después, otra con las coordenadas publicadas. El formato del edicto no dependía de quién lo firmaba.

La lista no es completa: el buscador ofrece más de quinientos avisos de toda la provincia para estos años, y sólo se revisó una parte de sus páginas. \pendiente{relación entre la comisión de la Resolución 122/16 y la determinación «por Administración» de la 389/16, ambas sobre el Quintín en la matrícula 4953; determinación resultante de la Resolución 290/16 (Chaile, matrícula 1534); coordenadas de la Resolución 279/17; expediente correcto de la Resolución 131/16; comisión de la Resolución 282/16; y revisión de las páginas restantes del listado de 2015--2017}
\end{cautela}

\subsection*{El arroyo Chaile, parcela por parcela (2016--2024)}

Las cuatro parcelas de la manzana 148 no fueron un caso aislado. Una búsqueda por el nombre del arroyo devuelve una serie que llega a 2024.

{\small
\setlength{\LTpre}{0.6em}\setlength{\LTpost}{1.2em}
\begin{longtable}{@{}p{2.4cm}p{3.3cm}p{6.2cm}@{}}
\toprule
\textbf{Acto} & \textbf{Expediente} & \textbf{Objeto y publicación} \\
\midrule
\endhead
Res.\ 251/16 (11/08/16) & 34-103.747/16 & \textbf{Comisión}: «arroyo Seco», matrícula 2821 (sección A, manzana 75, parcela 8D) y parte de la 4953. Torres y Caro. Edicto de junio de 2018: \textbf{publicada casi dos años después} \\
Res.\ 290/16 (07/09/16) & 34-178.741/16 & \textbf{Comisión}: matrícula 1534. La búsqueda por el expediente no devuelve determinación alguna \\
Res.\ 424 a 426/16 y 205 a 212/17 & cuatro expedientes & Manzana 148, matrículas 2342 a 2345. Con coordenadas \\
Sin número (jul.\ 2019) & 0090034-142.128/19 & \textbf{Comisión}: matrícula 5928. Torres y el ing.\ Julio César Zamora \\
Res.\ 327/19 (31/10/19) & 0090034-142.128/19 & \textbf{Determinación}: matrícula 5928. B.O.\ 3 y 4/12/2019 \\
Res.\ 75/2020 (03/03/20) & 0090034-103.747/2016 & \textbf{Determinación de línea de ribera y zona inundable}: «arroyos Vaqueros y Chaile», matrícula 2821 y parte de la 4953. B.O.\ 7 y 8/07/2020 \\
Res.\ 182/2021 (07/09/21) & 0090034-111.930/2021 & \textbf{Determinación}: catastro 1750. B.O.\ 22 y 23/09/2021 \\
Res.\ 51/2022 (27/04/22) & 0090034-170.889/2021 & \textbf{Determinación}: catastro 2318. B.O.\ 6 y 7/06/2022 \\
Res.\ 167/2022 (08/09/22) & 0090034-17.041/2022 & \textbf{Determinación}: catastro 2320. B.O.\ 13 y 14/10/2022 \\
Res.\ 206/2022 (04/11/22) & 0090034-177.288/2022 & \textbf{Determinación}: catastro 1095. B.O.\ 15 y 16/11/2022 \\
Res.\ 209/2022 (04/11/22) & 0090034-177.771/2022 & \textbf{Determinación}: catastro 1000. B.O.\ 24 y 25/11/2022 \\
Res.\ 258/2022 (29/12/22) & 0090034-216.166/2022 & \textbf{Determinación}: catastro 3604. B.O.\ 14 y 15/02/2023 \\
Res.\ 61/2024 (15/05/24) & 0090034-238.065/2023 & \textbf{Determinación}: colindante con la matrícula 2340. B.O.\ 14 y 18/06/2024 \\
\bottomrule
\end{longtable}
}

Desde 2019 ninguna de estas determinaciones publica coordenadas: todas remiten los «puntos de vinculación» a la sede. Las de 2019 y 2020 citan bien el Decreto 1989/02; las de 2021 a 2024 vuelven al «1898/02». Desde 2021 los edictos admiten el recurso por correo electrónico. Los firma la jefa del Programa Jurídico, Silvia Santamaría, sola o con la asesora legal. El costo de cada publicación pasó de doscientos pesos en 2016 a cuatro mil seiscientos veinte en 2024.

\begin{cautela}
Lo que esta serie muestra es un cambio de escala en la forma de fijar la ribera. Entre 2013 y 2015 las determinaciones del departamento recaían, en su mayoría, sobre fincas o fracciones rurales. Sobre el arroyo Chaile, a partir de 2016, recaen sobre parcelas: una manzana urbana de Vaqueros partida en cuatro, y después, entre 2019 y 2024, una determinación por catastro, a razón de una por año hasta 2021 y cinco en 2022. Es el ritmo al que un curso de agua queda rodeado de lotes. Este libro no identificó dónde están esos catastros ni quién pidió cada línea; su secuencia, en todo caso, coincide con la expansión de Vaqueros que documenta el capítulo \ref{cap:vaqueros}.

Dos rasgos merecen registro aparte. La Resolución 75/2020 es la única de la serie del Chaile que determina además una \textbf{zona inundable} ---en el departamento lo hicieron también la 348/13, la 14/16 y la 44/2021---, y lo hace sobre la confluencia con el arroyo Vaqueros; sus coordenadas, como las de las demás, no se publicaron. Y su comisión fue designada en 2016 para el «arroyo Seco», se publicó en 2018 y determinó en 2020 sobre los arroyos «Vaqueros y Chaile»: el mismo expediente nombra tres cursos distintos en cuatro años, del mismo modo que el de la Resolución 14/16 nombraba el Trancas y el Wierna. La matrícula 4953 de Finca Vaqueros aparece, con ésta, en cuatro expedientes distintos entre 2015 y 2020. \pendiente{titulares de los catastros 1000, 1095, 1750, 2318, 2320, 3604 y 5928 y de las matrículas 1534, 2340 y 2821 (ubicados en la capa catastral, salvo el 3604); zona inundable de la Resolución 75/2020; y relación de estas determinaciones con loteos o construcciones en Vaqueros}
\end{cautela}

\subsection*{El río Wierna, 2015--2024}

La misma búsqueda por nombre de curso, hecha sobre el Wierna, devuelve una serie paralela a la del Chaile.

{\small
\setlength{\LTpre}{0.6em}\setlength{\LTpost}{1.2em}
\begin{longtable}{@{}p{2.4cm}p{3.3cm}p{6.2cm}@{}}
\toprule
\textbf{Acto} & \textbf{Expediente} & \textbf{Objeto y publicación} \\
\midrule
\endhead
Res.\ 319/15 y 360/16 & 34-171.172/15 & Comisión y determinación: siete matrículas, entre ellas la \textbf{4097} y la 3438. Coordenadas remitidas a la sede \\
Res.\ 154/16 (17/05/16) & 34-259.760/14 & Determinación: matrícula 3722. Coordenadas al metro; publicada en noviembre de 2017 \\
Res.\ 304/17 (22/09/17) & 34-135.676/17 & \textbf{Comisión}: «río Las Nieves--Wierna», matrículas \textbf{4097} y 3438. Caro y Torres. Edicto del 15/05/2018, B.O.\ 8 y 11/06/2018 \\
Sin número (ago.\ 2019) & 0090034-117.269/2019 & \textbf{Comisión}: catastro 1821, \textbf{estancia Huerta Vieja}. Torres y el ing.\ Julio César Zamora. B.O.\ 23 y 26/08/2019 \\
Res.\ 181/2020 (09/12/20) & 0090034-117.269/2019 & \textbf{Determinación}: matrícula 1821. B.O.\ 3 y 4/03/2021 \\
Res.\ 44/2021 (25/02/21) & 0090034-208.131/2020 & \textbf{Determinación de línea de ribera y zona inundable}: matrícula \textbf{4097}. B.O.\ 5 y 6/04/2021 \\
Res.\ 174/2021 (24/08/21) & 0090034-57.407/2021 & \textbf{Determinación}: catastro 47. B.O.\ 30 y 31/08/2021 \\
Res.\ 94/2022 (21/06/22) & 0090034-70.441/2022 & \textbf{Determinación}: matrículas \textbf{4097} y 6663. B.O.\ 19 y 20/07/2022 \\
Res.\ 23/2024 (08/03/24) & 0090034-81.638/2022 & \textbf{Determinación}: catastro 3714 «y parcelas allí citadas». Autorizado a diligenciar: el ing.\ civil Diego Dean. Edicto del 05/11/2024, B.O.\ 26 y 27/12/2024 \\
\bottomrule
\end{longtable}
}

\begin{cautela}
Una matrícula se repite. La \textbf{4097} aparece en la determinación de 2016, en la comisión de 2017 sobre el «Las Nieves--Wierna», en la determinación con zona inundable de 2021 y en la de 2022, esta vez junto a una matrícula de numeración mucho más alta, la 6663. Cuatro actos sobre el mismo inmueble en seis años no son un relevamiento: son las etapas de algo que ocurre en él. La aparición de una matrícula nueva a su lado sugiere una subdivisión, pero eso sólo lo diría el folio real. Es, con la Resolución 75/2020 del Chaile, una de las dos determinaciones del departamento con \textbf{zona inundable} que no publican sus puntos ---la 348/13 y la 14/16 sí los publicaron---.

El nombre «Las Nieves--Wierna» de la Resolución 304/17 importa por otra razón: el Las Nieves es el curso que la comunidad Kondorwaira nombra como su cuenca de habitación y el que el Estado dejó de aforar en 1961 (capítulo \ref{cap:resistencias}). La comisión que en 2017 va a fijar su ribera se publica ocho meses después de designada.

Las dos comisiones de 2019 ---la de la estancia Huerta Vieja sobre el Wierna y la de la matrícula 5928 sobre el Chaile--- se publicaron con un mes de diferencia, con los mismos dos técnicos y con una fórmula que ninguna otra usa: «se determinó la Comisión Técnica», sin número ni fecha de resolución. El edicto no dice qué acto la designó.

Y la Resolución 428/15, que este capítulo sumó a la serie de 2013 a 2015, reúne en un solo edicto tres de los defectos que la serie venía mostrando por separado: el número de expediente difiere en un dígito del que el índice asigna al mismo inmueble (34-185.548/14 contra 34-485.548/14), el curso es otro (río La Caldera contra arroyo Isasmendi) y el departamento está mal escrito (La Candelaria). Publica, en cambio, diecisiete coordenadas. \pendiente{folio real de las matrículas 4097 y 6663; zona inundable de la Resolución 44/2021;  aclaración del número y del curso del expediente de la matrícula 2020}
\end{cautela}

\subsection*{El río y el arroyo Vaqueros, 2016--2025}

{\small
\setlength{\LTpre}{0.6em}\setlength{\LTpost}{1.2em}
\begin{longtable}{@{}p{2.4cm}p{3.3cm}p{6.2cm}@{}}
\toprule
\textbf{Acto} & \textbf{Expediente} & \textbf{Objeto y publicación} \\
\midrule
\endhead
Res.\ 009/12 & --- & Determinación: río Vaqueros, colindante con la matrícula \textbf{216}. Conocida sólo por su rectificación de 2025 \\
Res.\ 147/16, 414/16 y 144/17 & 34-202.169/15 & Comisión, determinación y ampliación: río Vaqueros, matrícula 40 (véase arriba) \\
Res.\ 338/16 (11/10/16) & 34-164.459/16 & \textbf{Comisión}: río Vaqueros, matrícula 200. Torres y Caro. Edicto del 20/04/2018, \textbf{publicado el 24 y 25/01/2019}, veintisiete meses después de la resolución \\
Res.\ 23/2021 (28/01/21) & 0090034-203.627/2020 & \textbf{Determinación}: arroyo Vaqueros, matrícula 1566. B.O.\ 9 y 10/02/2021 \\
Res.\ 125/2021 (14/06/21) & 0090034-84.440/2021 & \textbf{Determinación}: arroyo Vaqueros, catastro \textbf{3017}. B.O.\ 25 y 28/06/2021 \\
Res.\ 107/2022 (07/07/22) & 0090034-88.235/2022 & \textbf{Determinación}: arroyo Vaqueros, matrícula \textbf{3017}, otra vez. B.O.\ 12 y 13/07/2022 \\
Res.\ 56/2025 (20/03/25) & 0090034-185.051/2024 & Determinación: arroyo Vaqueros, catastro 1122 (véase el comienzo de este capítulo) \\
Res.\ 57/2025 (20/03/25) & 0090034-203.058/2024 & \textbf{Determinación}: arroyo Vaqueros, matrícula \textbf{216}. Autorizado a diligenciar: el geólogo Víctor García. B.O.\ 16 y 21/04/2025 \\
Res.\ 130/2025 (09/06/25) & 0090034-69.337/2025 & \textbf{Rectificación parcial de la Resolución 009/12}: río Vaqueros, matrícula \textbf{216}. Autorizado a diligenciar: el licenciado Enrique Chalabe. B.O.\ 11 y 14/07/2025 \\
\bottomrule
\end{longtable}
}

\begin{cautela}
El patrón del Chaile y del Wierna se repite. Desde 2019 ninguna determinación publica sus puntos, y dos inmuebles vuelven: la matrícula \textbf{3017} tiene determinación en 2021 y otra en 2022, y la \textbf{216} tiene una línea de 2012, una determinación nueva en marzo de 2025 y, tres meses después, la rectificación parcial de la de 2012. Tres actos sobre el mismo predio, dos de ellos el mismo año, con dos profesionales distintos autorizados a gestionar su publicación. Uno de ellos, el licenciado Chalabe, había integrado en 2016 la comisión del arroyo Chaile: como el ingeniero Forns en el Quintín, aparece primero en una comisión del Estado y después como gestor de un trámite ante el mismo organismo. Los edictos no dicen por cuenta de quién.

La comisión de la matrícula 200 establece, además, la demora más larga de toda la serie: veintisiete meses entre la resolución y su publicación. Durante ese lapso rigió la prohibición de modificar las márgenes que la propia resolución imponía, sin que los ribereños pudieran conocerla por el Boletín. \pendiente{texto de la Resolución 009/12 y de su rectificación de 2025; y relación entre las dos determinaciones de la matrícula 3017}
\end{cautela}

\section{Lo que dice la capa catastral}

Los edictos nombran inmuebles por su número; la capa «Catastro Urbano y Rural» de la Dirección General de Inmuebles, descargada del Geoportal de IDESA en septiembre de 2026, permite ubicarlos y medirlos. El organismo advierte que esa cartografía es orientativa, un croquis de ubicación y no un plano aprobado, y así se la usa aquí. El departamento tiene en ella 6.219 parcelas: 5.800 urbanas, 225 subrurales y 194 rurales. La lámina \ref{fig:catastro} las ubica.

\begin{figure}[p]
\centering
\includegraphics[width=\textwidth,height=0.84\textheight,keepaspectratio]{fig-catastro-ribera.png}
\caption[Las parcelas de las líneas de ribera y de las concesiones de agua]{\textbf{Las parcelas de las líneas de ribera y de las concesiones de agua.} A: el departamento, con las parcelas sobre las que recayó una determinación o comisión de línea de ribera (rojo), una concesión de agua (azul) o una asignación de riego por subdivisión (verde), y los 247 lotes y el remanente del plano de El Durazno (violeta). B: Vaqueros, donde los círculos marcan las parcelas menores de dos hectáreas; en trazo oscuro, las líneas de ribera cuyas coordenadas publicó el Boletín (Resoluciones 80/14, 201/13, 414/16, 144/17, 14/16 y 205 a 212/17), y en punteado celeste la línea de inundación de la 14/16. Los blancos dentro del departamento son sectores sin parcela en la capa. Elaboración propia con la capa «Catastro Urbano y Rural de la Provincia de Salta» de la Dirección General de Inmuebles (Geoportal IDESA, descargada en septiembre de 2026; cartografía orientativa según el propio organismo), la cartografía de radios del Censo 2022 (INDEC) y las coordenadas publicadas en los edictos. Coordenadas convertidas desde POSGAR 94, faja 3. Elaboración propia. \textbf{El script que dibuja esta lámina no está todavía en el repositorio} \texttt{mapas\_caldera}: \texttt{scripts/mapa\_ribera.py} dibuja la lámina \ref{fig:riberaedictos}, no ésta (apéndice \ref{cap:fuentes}).}\label{fig:catastro}
\end{figure}


{\small
\setlength{\LTpre}{0.6em}\setlength{\LTpost}{1.2em}
\begin{longtable}{@{}p{1.5cm}p{3.7cm}p{1.9cm}p{4.8cm}@{}}
\toprule
\textbf{Catastro} & \textbf{Acto del libro} & \textbf{Superficie} & \textbf{Categoría y finca} \\
\midrule
\endhead
\multicolumn{4}{@{}l}{\textit{Líneas de ribera sobre lotes urbanos}} \\
216 & Res.\ 009/12, 57/2025 y 130/2025 & 0,14 ha & Urbano, Vaqueros \\
3017 & Res.\ 125/2021 y 107/2022 & 0,38 ha & Urbano, Vaqueros \\
1566 & Res.\ 23/2021 & 0,64 ha & Urbano, Vaqueros \\
1122 & Res.\ 56/2025 & 0,58 ha & Urbano, Vaqueros \\
2342--2345 & Res.\ 205 y 210--212/17 & 0,08--0,09 ha c/u & Urbano, planos 362 y 1230 \\
2318 y 2320 & Res.\ 51 y 167/2022 & 0,05 ha c/u & Urbano, Vaqueros \\
1000, 1095, 1750, 2340 & Res.\ 209 y 206/2022, 182/2021, 61/2024 & 0,10--0,96 ha & Urbano, Vaqueros \\
2821 & Res.\ 75/2020 (con zona inundable) & 0,24 ha & Urbano, Vaqueros \\
4027 & Res.\ 94/12 (advertencia de no construir) & 1,35 ha & \textbf{Urbano} \\
3756 & Res.\ 292/13 y 142/2025 & 23,6 ha & \textbf{Urbano} (en 2013, «fracción 48») \\
4146 & Res.\ 291/16, 279/17, 108/2025 & 5,05 ha & Urbano, La Calderilla \\
\multicolumn{4}{@{}l}{\textit{Líneas de ribera sobre fincas}} \\
40 & Res.\ 147/16, 414/16, 144/17 & 2,94 ha & Urbano, sección B \\
200 & Res.\ 338/16 & 4,61 ha & Subrural, finca Cerro de Buena Vista \\
23 & Res.\ 198/15 y 322/15 & 5,64 ha & Subrural, finca Vaqueros \\
4097 & Res.\ 360/16, 304/17, 44/2021, 94/2022 & \textbf{901 ha} & Rural, finca Vaqueros (plano 860) \\
3722 & Res.\ 74/15, 293/15, 154/16 & 320 ha & Rural, finca Wierna \\
2020 & Res.\ 428/15 & 39 ha & Rural, finca Wierna \\
1821 & Res.\ 181/2020 & 215 ha & Rural, finca Huerta Vieja \\
47 & Res.\ 174/2021 & 380 ha & Rural, La Mesada, fracción B \\
3438 & Res.\ 319/15, 360/16, 304/17 & 78 ha & Subrural \\
5928 & Res.\ 327/19 & 30 ha & Rural, finca Vaqueros \\
1499 & Res.\ 282/16 & 51 ha & Rural, finca San Jorge \\
\multicolumn{4}{@{}l}{\textit{Concesiones de agua}} \\
4095 & Las Vertientes (2017) & 7,0 ha & Subrural, finca Vaqueros (plano 860, lindero de la 4097) \\
1782 & Skru S.A.\ (2019) & 94 ha & Rural, finca Campo Alegre \\
5925 & El Chalchal S.R.L.\ y otro (2025) & 573 ha & Rural, finca Chaguadero y San José \\
3307 & Valle Hermoso S.R.L.\ (2019) & 10,9 ha & Urbano \\
1673 & Arzobispado de Salta (2024) & 104 ha & Rural, finca Orco Huasi \\
6582 & Riego del catastro 2907 (2021) & 32 ha & Rural, finca \textbf{Sausal o Curuzú} \\
5009 y 5010 & Riego del catastro 3721 (2018, 2022) & 62 y 238 ha & Rurales, finca Wierna, linderas \\
3436 y 3437 & Riego del catastro 169 (2022) & 182 y 169 ha & Rurales, fincas Loma Pelada y Antilla, linderas \\
7155 y 7156 & Riego de los catastros 84, 2264 y 2932 (2023) & 2,0 y 10,8 ha & Subrurales, finca La Calderilla \\
3127 & Riego del catastro 38 (2005) & 49 ha & Rural, finca San Antonio \\
2762 & Vivero (2023) & 0,81 ha & Urbano \\
\bottomrule
\end{longtable}
}

\begin{cautela}
La capa confirma lo que los edictos sugerían y le pone medida. Las líneas de ribera del arroyo Vaqueros y del Chaile posteriores a 2019 no recaen sobre fincas sino sobre \textbf{lotes urbanos de entre quinientos metros cuadrados y una hectárea}. El inmueble de la matrícula 216, con una línea de 2012, otra de marzo de 2025 y la rectificación de junio, mide mil cuatrocientos metros cuadrados. La determinación de la ribera, que en 2013 se hacía para fracciones rurales, es hoy un trámite de lote.

Tres casos merecen registro aparte. La matrícula 4027, cuyo propietario fue advertido por escrito en 2012 de que no construyera viviendas cerca del río Castellanos, figura hoy como parcela \textbf{urbana}. La 3756, que en 2013 era la «fracción 48» de Finca Vaqueros, figura como urbana de veintitrés hectáreas cuando en 2025 se aprueba su línea. Y la 6582, a la que en 2021 se pasó el riego de la 2907, es la finca \textbf{«Sausal o Curuzú»}, la misma que el juez mandó deslindar en 1921 junto con Cerro de Buena Vista (capítulo \ref{cap:loteo}); la matrícula 200, con comisión de 2016 sobre el río Vaqueros, lleva el nombre de esta última.

La capa también muestra lo que ya no está. No figuran las matrículas de origen de los loteos y de las subdivisiones ---3616 (La Misión III), 3709 (Chacras del Gallinato), 3989 (Las Vertientes), 1905 y 1713 (Cerro de Buena Vista), 4366 y 6558 (El Durazno), 4953 (Finca Vaqueros), 3721, 2907, 126, 169 y 38---: se dieron de baja al dividirse. Las cadenas que los edictos de riego describen se verifican en el mapa: las matrículas 5009 y 5010 son linderas y salen de la 3721; las 3436 y 3437, de la 169; la 6582 sustituye a la 2907. La 4095 de Las Vertientes y la 4097 comparten el plano 860 y lindan. Una sola discordancia: el catastro 2264, que el edicto de 2023 da por anulado, sigue en la capa.

Por último, una conjetura que la capa descarta: los catastros 5925 y 5928, de numeración casi contigua, no son vecinos; los separan novecientos metros y dos planos distintos.
\end{cautela}

\section{El artículo 172, que ninguna fuente anterior había traído}

El edicto de 2013 transcribe el texto de una norma que cambia el peso de todo este capítulo.

\begin{ficha}{Art.\ 172 del Código de Aguas (Ley 7017) --- texto transcripto en el edicto}
``Zonas inundables. La Autoridad de Aplicación, para las zonas ribereñas determinará los sectores que puedan ser afectados por inundaciones, \textbf{en los cuales no se permitirá la instalación de asentamientos poblacionales} o fijación de obstáculos que puedan afectar el curso de las aguas. \textbf{Las nuevas construcciones o plantaciones que se efectúen en zonas ribereñas deberán ser autorizadas previamente por la Autoridad de Aplicación}, debiéndose proyectar y ejecutar las obras necesarias para prevenir el riesgo de inundación.''
\end{ficha}

\begin{cautela}
Léase con cuidado, porque es la norma más exigente que este libro encontró y conviene no forzarla.

Dice dos cosas distintas. La primera es una \textbf{prohibición}: en los sectores que la Autoridad determine como afectables por inundaciones no se permite la instalación de asentamientos poblacionales. La segunda es un \textbf{requisito de autorización previa}: toda nueva construcción en zona ribereña necesita permiso de la Autoridad de Aplicación, y deben proyectarse y ejecutarse las obras que prevengan el riesgo.

Ambas dependen de que exista determinación previa. La prohibición opera sobre ``los sectores que'' la Autoridad haya determinado; el permiso opera sobre ``zonas ribereñas'', que la línea de ribera define. En el río La Caldera, matrícula 1590, esa determinación existe desde noviembre de 2013 e incluye la zona inundable. En los arroyos El Durazno y Guaranguay, matrícula 4366, existe línea de ribera desde diciembre de 2014; si incluyó zona inundable es lo que no se sabe.

De ahí las dos preguntas que este capítulo debe formular y no responder. \textbf{Primera:} ¿la Resolución 383/14 determinó zona inundable sobre los arroyos, y qué superficie del loteo El Durazno queda dentro? \textbf{Segunda:} las construcciones levantadas en zona ribereña de esos arroyos, ¿contaron con la autorización previa que exige el artículo 172, y con las obras de prevención que la norma manda proyectar y ejecutar?

Nada permite afirmar que no las hayan tenido. Afirma que la norma las exigía, que el organismo competente es el mismo que intervino a fojas 59 del expediente del loteo, y que la respuesta está en ese expediente y en el de la determinación. No es una hipótesis: es un permiso que se otorgó o no se otorgó, y consta o no consta.
\end{cautela}

Una precisión sobre el error de numeración, que este edicto obliga a reformular. El de 2013 también cita ``Dcto.\ Nº 1898/02, art.\ 10'', igual que el de julio de 2014 ---el de la Resolución 185/14--- y el de 2015. El de 2019 y el de julio de 2025 citan 1989/02, que es el correcto; el de abril de 2025 sobre el arroyo Vaqueros vuelve a decir 1898/02. Y antes de la serie, el edicto de mayo de 2012 sobre Chacras del Gallinato, firmado por otro funcionario, también cita correctamente el 1989/02: la transposición no viene del origen del régimen.

No fue entonces un desliz aislado: fue un \textbf{error de plantilla que atraviesa toda la serie}. Las tablas de este mismo capítulo lo registran en 2013, 2014 y 2015, en las Resoluciones 282/16, 389/16 y 279/17 y en la serie del arroyo Chaile de 2021 a 2024; el edicto de 2019 sobre el río Bermejo y el de julio de 2025 citan el número correcto, y el de abril de 2025 vuelve al error. Las dos versiones conviven desde 2013. Es el segundo error de plantilla que este libro documenta --- el primero es el municipio de Rivadavia Banda Sur en el Plan de Desarrollo Urbano Ambiental --- y ambos comparten el mismo mecanismo: un texto que se reutiliza y que nadie relee.

La pregunta que este capítulo había dejado abierta --- si la Ley 7017 fija, como su análoga cordobesa, un plazo para levantar los planos de zonas inundables --- ya tiene respuesta. Obtenido el texto completo del Código de Aguas y de su reglamentación: \textbf{no lo fija}.

El artículo 194 del código cordobés, de 1973, ordena a la Autoridad levantar esos planos ``dentro de los diez años''. El artículo 172 salteño, de 1998, reproduce el resto del contenido --- la prohibición de asentamientos, la autorización previa de construcciones --- y \textbf{omite el plazo}. La obligación de determinar los sectores inundables existe; el término para cumplirla, no.

Es una diferencia de una línea y tiene consecuencia práctica: sin plazo no hay incumplimiento de plazo. La determinación de zonas inundables en el departamento La Caldera podía hacerse en 1999 o en 2030 sin que ninguna fecha se venciera.

\section{Y el certificado no es un trámite: es un recaudo indispensable}

La reglamentación del Código de Aguas contiene, en el artículo que reglamenta el 172, la norma que la presentación del COHIFE mencionaba sin citar.

\begin{ficha}{Art.\ 172º de la reglamentación --- texto}
``La Autoridad de Aplicación es el organismo competente para expedir las certificaciones de \textbf{`no inundabilidad'}, que constituyen \textbf{un recaudo indispensable para todo asentamiento poblacional, urbanización, o plantación en zona ribereña}. Ello, previo cobro a los interesados de un arancel, que fijará para cada caso la Autoridad de Aplicación.''
\end{ficha}

``Recaudo indispensable''. No es una constancia recomendable ni un requisito que la autoridad pueda dispensar: la norma lo califica de indispensable, y lo hace \textbf{para todo asentamiento poblacional o urbanización en zona ribereña}, sin distinguir tamaño, época ni jurisdicción.

El loteo El Durazno es una urbanización. Está atravesado por dos arroyos con línea de ribera determinada desde diciembre de 2014. El organismo competente para expedir el certificado es el mismo que intervino a fojas 59 del expediente de aprobación de su plano, en 2021.

La pregunta que este capítulo formuló hace unas páginas queda ahora apoyada en el texto de una norma y no en una presentación: \textbf{¿se expidió el certificado de no inundabilidad?} Y como el propio artículo prevé el cobro de un arancel, hay además un rastro contable: un certificado expedido dejó un pago.

\section{Dos artículos vecinos que el capítulo necesitaba}

\begin{ficha}{Arts.\ 171 y 173 de la Ley 7017 --- texto}
\textbf{Art.\ 171 --- Información previa.} ``Previo al otorgamiento, por parte de la autoridad competente, de permisos o concesiones de uso de cauces, márgenes o \textbf{extracción de áridos}, la Autoridad de Aplicación evaluará si el permiso o concesión afectará desfavorablemente las riberas o el flujo de aguas, si así fuera desaconsejará el otorgamiento del permiso o concesión, o se exigirá la construcción de las obras necesarias y la aplicación de las técnicas de operación adecuadas para prevenir los peligros y los daños.''

\textbf{Art.\ 173 --- Penalidades.} Las infracciones al artículo 172 se sancionan con multa y sanciones conminatorias. ``Sin perjuicio de ello la Autoridad de Aplicación podrá \textbf{ordenar la demolición de las obras o la destrucción de los obstáculos} o demolerlos o destruirlos por cuenta del infractor.''
\end{ficha}

El artículo 171 conecta con lo que, según la prensa, se ordenó reparar en la ejecución en noviembre de 2025 ---acto que la lista de actuaciones consultada no registra---: el daño causado por la extracción irregular de áridos. La norma exigía, desde antes, una evaluación previa de impacto sobre riberas y flujo para cada permiso de extracción. Si esas evaluaciones existen, están en la Autoridad de Aplicación.

El artículo 173 es el que más pesa, y conviene enunciarlo con cuidado. \textbf{La provincia tiene, desde 1998, la potestad de ordenar la demolición de obras levantadas en violación del artículo 172.} No es una facultad que este libro reclame ni recomiende --- se trata de viviendas habitadas por familias, y ninguna solución que pase por demolerlas sería una solución. Lo que corresponde registrar es de otro orden: la herramienta jurídica más severa del régimen hídrico provincial existía y no fue necesario discutirla, porque nunca se llegó al momento en que hubiera que decidir si aplicarla.

Cuándo se estuvo cerca de ese momento, y quién no lo advirtió, es la pregunta que atraviesa este capítulo desde su primera página.

\section{Hay rectificaciones en el departamento, y son recientes}

El capítulo \ref{cap:vaqueros} planteaba, sobre la base del mismo material, que existe un mecanismo de modificación de la línea de ribera a pedido de desarrolladores. La existencia del mecanismo quedó confirmada por el artículo 1º del Anexo I --- la determinación procede a solicitud de parte --- y por el artículo 126 in fine del Código de Aguas. Y ahora hay casos fechados en el propio departamento. El primero que se transcribe no es, en rigor, una rectificación sino una determinación demorada casi doce años; las rectificaciones son las que le siguen.

\begin{ficha}{Resolución Nº 000142/2025 --- B.O.\ Nº 22.007 y 22.008, 8 y 11 de agosto de 2025}
Expediente 0090034-160869/2013-0/4. Aprueba, el 30 de junio de 2025, la determinación de la línea de ribera realizada por la Comisión Técnica en ambas márgenes del \textbf{arroyo Quintín}, a la latitud del \textbf{catastro 3756}. Puntos de vinculación a disposición en la sede. Autorizado a diligenciar la publicación: \textbf{Ing.\ Antonio Alejandro Forns}\index{Forns, Antonio Alejandro}. Firma la jefa del Programa Jurídico, Silvia F.\ Santamaría. Cita correctamente el Decreto 1989/02 y admite el recurso por correo electrónico. Edicto del 7 de julio de 2025.
\end{ficha}

\begin{cautela}
Es la determinación que cierra la comisión de la Resolución 292/13: el mismo expediente, el mismo catastro y el mismo arroyo, \textbf{casi doce años después}. En el camino el número del expediente migró al formato nuevo y cambiaron el asesor y el formato del edicto. Lo que no cambió es un nombre: el ingeniero Forns, que integró la comisión de 2013, es en 2025 quien está autorizado a diligenciar la publicación. El edicto no dice en representación de quién. En la Resolución 108/2025 la misma función la cumple otro ingeniero. Qué determinó la comisión en 2013 y qué se aprobó en 2025 ---si es la misma línea o una actualizada por la migración del cauce--- es lo que diría el expediente. \pendiente{expediente 0090034-160869/2013, con el informe de la comisión y los fundamentos de la Resolución 142/2025}
\end{cautela}

\begin{ficha}{Resolución Nº 000108/2025 --- B.O.\ Nº 21.962, 30/05 y 02/06/2025}
Expediente Nº 0090034-242142/2024-0. La Secretaría de Recursos Hídricos hace saber que por Resolución Nº 000108/2025 del \textbf{9 de mayo de 2025} \textbf{se rectificó en su parte pertinente la determinación de la línea de ribera aprobada por las Resoluciones Nº 279/17 y 306/19}, realizada por la Comisión Técnica en ambas márgenes del \textbf{río La Caldera}, en la zona de la \textbf{Matrícula Nº 4146} del departamento La Caldera.

Puntos de vinculación a disposición en la sede. Autorizado a diligenciar: \textbf{Ing.\ Jesús E.\ Flores}. Firma: Dra.\ Silvia F.\ Santamaría\index{Santamaría, Silvia F.}, Jefa de Programa Jurídico. Apelable ante el Tribunal de Aguas (Decreto 1989/02, art.\ 10) en diez días hábiles administrativos.
\end{ficha}

\begin{cautela}
Conviene separar tres asuntos.

\textbf{Primera: la serie del río La Caldera es larga.} La comisión se designó por Resolución 291/16 (véase el apartado sobre 2016 y 2017); la línea de ribera se determinó por Resolución 279/17, se volvió a tratar por Resolución 306/19 y se rectificó por Resolución 108/2025. Tres actos en ocho años sobre el mismo curso, en un departamento donde entre medio ocurrieron dos aluviones y una condena judicial.

\textbf{Segunda: aparece una matrícula nueva.} La 4146. El libro venía siguiendo la 4366 --- zona de la determinación de 2014 sobre los arroyos --- y la 6.558 --- inmueble del loteo El Durazno. Tres matrículas del mismo departamento, sobre las que el Estado fijó o rectificó líneas de ribera y aprobó un fraccionamiento. Establecer su relación es, otra vez, el antecedente dominial que falta.

\textbf{Tercera: qué se rectificó y por qué.} El edicto dice ``en su parte pertinente'' y no dice más. Puede tratarse de una corrección técnica menor, de una actualización por migración del cauce --- que el artículo 7 del Anexo I contempla expresamente --- o de una modificación de alcance. Este libro \textbf{no lo sabe} y no debe elegir entre las tres. Lo que corresponde es pedir el expediente 0090034-242142/2024-0, que es corto y está identificado.
\end{cautela}

Queda resuelta, además, la discrepancia que este capítulo había señalado entre los decretos 1989/02 y 1898/02. El correcto es \textbf{1989/02}: así lo cita la presentación oficial, así lo cita el edicto del río Bermejo y así lo cita el edicto de julio de 2025; el de abril de ese año sobre el arroyo Vaqueros, en cambio, repite el error. La mención a ``1898/02'' en el edicto de la Resolución 383/14 es una transposición de dígitos \textbf{en el propio edicto}.

Es un error menor y sin consecuencia jurídica, pero se repite en edictos de 2013 a 2017 y de 2021 a 2025, falta en otros de esos mismos años ---y en los de 2019 y julio de 2025---, y reaparece en el de abril de 2025: es un defecto de plantilla que convive con la versión corregida. Conviene registrarlo por lo que dice del control de calidad de las publicaciones: durante años, el instrumento que informaba a los vecinos bajo qué norma podían apelar citó mal el número de esa norma.

Se obtuvo el texto del artículo 126 y un edicto de 2019 del mismo organismo que permite ver el procedimiento completo. Ambos corrigen datos que este apartado había tomado de un resumen intermedio, antes de identificar la presentación de Torres.

\begin{ficha}{Art.\ 126 del Código de Aguas (Ley 7017) --- texto}
La Autoridad de Aplicación determina la línea de los cursos naturales \textbf{conforme al criterio del artículo 2.577 del Código Civil} y según el procedimiento técnico que fije la reglamentación, \textbf{dando intervención en la operación a los interesados}. Las cuotas determinantes se anotan en el Libro de Catastro que prevé el Código. Y en su última oración: \textbf{``La Autoridad de Aplicación podrá rectificar la línea de ribera cuando el cambio de las circunstancias lo haga necesario.''}
\end{ficha}

\begin{cautela}
Dos observaciones sobre este artículo.

La primera es una cuestión abierta que este libro señala sin resolver. El artículo remite al \textbf{artículo 2.577 del Código Civil} de Vélez, cuerpo derogado en agosto de 2015 y reemplazado por el Código Civil y Comercial, cuyo artículo 235 delimita el río por la línea de ribera que fija \textbf{el promedio de las máximas crecidas ordinarias}. Qué criterio rige hoy en Salta --- si la remisión debe entenderse dinámica hacia el texto vigente, o si subsiste el criterio del código derogado --- es una pregunta jurídica que excede a este libro y que corresponde formular a la Autoridad de Aplicación. La determinación de 2014 que aquí importa es anterior a esa derogación.

La segunda es que la \textbf{rectificación está prevista en la ley}. La presentación de Torres describe que proyectos privados solicitan corrimientos de la línea; el artículo 126 in fine muestra que la potestad de rectificar existe, aunque el texto la atribuye a la Autoridad ``cuando el cambio de las circunstancias lo haga necesario'' y no a pedido de un desarrollador. Que exista la puerta no prueba quién la usa ni para qué. Los expedientes de rectificación del departamento siguen siendo el documento que lo establecería.
\end{cautela}

\section{Once años después, el edicto sigue siendo el mismo}

El núcleo de este capítulo ---la serie de veinticinco actos--- ocurrió entre 2013 y 2015. Queda la pregunta de si aquella
forma de publicar era propia de esos años, de ese departamento o de aquel firmante. Un edicto de
2025 la contesta.

\begin{ficha}{Boletín Oficial Nº 21.987, 10 y 11 de julio de 2025}
Expediente 0090034-194687/2024-0. La Secretaría de Recursos Hídricos hace saber que por
\textbf{Resolución Nº 000111/2025} del 30 de junio de 2025 se aprobó la determinación de línea
de ribera en ambas márgenes del \textbf{río Ancho}, a la latitud del catastro 9.226 del
\textbf{departamento Cerrillos}, «\textit{cuyos puntos de vinculación \textbf{se encuentran a
disposición de los interesados en la sede} de la Secretaría de Recursos Hídricos}».

Firma la \textbf{Dra.\ Silvia F.\ Santamaría}\index{Santamaría, Silvia F.}, Jefa de Programa
Jurídico. Arancel: \textbf{once mil novecientos pesos}.
\end{ficha}

\textbf{Es la misma fórmula, palabra por palabra, que la Resolución 383/14.} Otro departamento,
otra cuenca, otro firmante, once años después.

La fecha de la resolución es la que da el edicto, y no encaja en la numeración de la
Secretaría: la 108/2025 es del 9 de mayo, la 130/2025 del 9 de junio y la 142/2025 del 30 de
junio, el mismo día que el edicto atribuye a la 111/2025. \pendiente{fecha de la Resolución 111/2025 en el registro de la Secretaría de Recursos Hídricos}

\begin{cautela}
Este hallazgo corta en dos direcciones y las dos importan.

\textbf{Debilita cualquier lectura particular.} No hubo una decisión tomada sobre El Durazno y
Guaranguay: hay \textbf{una forma de redactar edictos} que el organismo viene usando, que
alterna con la otra ---la que publica coordenadas--- y que sigue vigente. Este libro nunca
afirmó que la omisión fuera deliberada; ahora tiene además una razón positiva para no
afirmarlo.

\textbf{Y refuerza la tesis estructural.} Lo que el capítulo \ref{cap:opacidad} describe no es
una conducta sino un criterio disponible: dos plantillas, ninguna norma que mande elegir una, y
nadie obligado a explicar la elección. \textbf{Que once años después convivan las dos es la
prueba más limpia de que no hay criterio escrito.}

\textbf{Lo que este edicto no prueba}, por sí solo, es cómo se publicaron las determinaciones de La Caldera entre 2016 y 2024; las series del Chaile, del Wierna y del Vaqueros que este capítulo reúne muestran que desde 2019 ninguna publica sus coordenadas. En 2025 la fórmula sí se repitió en el departamento ---la Resolución 56/2025 sobre el arroyo Vaqueros y la 108/2025 sobre el río La Caldera remiten sus puntos a la sede---, pero el apéndice \ref{cap:pedidos} pide las resoluciones 279/17 y 306/19, que son las que faltan para cerrar la serie del río La Caldera.
\end{cautela}

\subsection*{Y en 2026, otra vez}

No es un caso aislado de 2025. Nueve meses después vuelve a ocurrir, en un tercer departamento.

\begin{ficha}{Boletín Oficial Nº 22.160, 1 y 6 de abril de 2026}
Expediente 0090034-130747/2025-0. La \textbf{Subsecretaría} de Recursos Hídricos hace saber que
por \textbf{Resolución Nº 000016/2026} del 4 de marzo de 2026 se aprobó la determinación de
línea de ribera en ambas márgenes del \textbf{río Calchaquí}, en zona de la matrícula 538 del
\textbf{departamento San Carlos}, «\textit{cuyos puntos de vinculación se encuentran a
disposición de los interesados en la sede de la Secretaría de Recursos Hídricos}».

Firma la \textbf{Dra.\ Silvia F.\ Santamaría}, ahora \textbf{Directora de la Dirección Legal y
Técnica}. Arancel: \textbf{diecisiete mil cuarenta pesos}.
\end{ficha}

\textbf{Tres departamentos, tres cuencas, tres años, la misma frase.} Cerrillos en 2025, San
Carlos en 2026, La Caldera en 2014. Lo que este capítulo documentó como un episodio es
\textbf{la forma habitual de un organismo}.

\begin{cautela}
\textbf{Y este edicto resuelve un pedido del apéndice \ref{cap:pedidos}.}

Ese apéndice preguntaba si el organismo hídrico había cambiado de rango entre 2024 y 2025. La
respuesta está en el encabezado de los dos edictos: en julio de 2025 firma la
\textbf{Secretaría} de Recursos Hídricos; en marzo de 2026, la \textbf{Subsecretaría}.
\textbf{El organismo bajó de rango entre ambas fechas.} La ventana se acorta: el Decreto 52/26, publicado el 13 de febrero de 2026, designa al Ing.\ Julio Mauricio \textbf{Romero Leal}\index{Romero Leal, Mauricio} \textbf{Subsecretario} de Recursos Hídricos. Es el mismo funcionario que en 2022 y 2023 figuraba como \textbf{secretario} del área (capítulo \ref{cap:amparo}): el organismo perdió rango y su titular volvió a él con el cargo nuevo. El cambio es parte del reordenamiento de fines de 2025. La \textbf{Ley 8.511} (B.O.\ del 26/11/2025) suprimió el Ministerio de Infraestructura, pasó las tierras fiscales a la Jefatura de Gabinete (capítulo \ref{cap:tierrafiscal}) y fijó un tope de treinta secretarías de Estado (art.\ 30). En la letra de la ley, el uso y aprovechamiento de los recursos hídricos corresponde al Ministerio de \textbf{Producción y Minería}, junto con el contralor de las explotaciones mineras, la autoridad ambiental y el ordenamiento territorial (art.\ 25, incs.\ 7, 9, 13 y 14). Pero la estructura de ese ministerio, aprobada por la Decisión Administrativa 565/25 con vigencia al 15 de diciembre de 2025 (B.O.\ Nº 22.100 del 30/12/2025; lámina \ref{lam:da565}), \textbf{no incluye al organismo hídrico}. Su organigrama tiene cuatro secretarías ---Ambiente y Desarrollo Sustentable, Industria y Comercio, Minería y Desarrollo Agropecuario--- y ninguna unidad de recursos hídricos en ningún nivel. La Subsecretaría quedó en la \textbf{Jefatura de Gabinete de Ministros}, a la que la misma ley asigna las obras de aprovechamiento de cuencas y recursos hídricos (art.\ 19, inc.\ 16): el Decreto 52/26, del 10 de febrero de 2026 (B.O.\ Nº 22.131), que designa a Romero Leal, lleva ese encabezado y lo refrenda el jefe de Gabinete. Lo que el organigrama sí muestra es que el ambiente conservó el rango de Secretaría y que ambiente y minería quedaron como secretarías hermanas bajo el mismo ministro. \textbf{El agua, en cambio, perdió rango y quedó fuera del ministerio que la ley nombra competente.} Quedó, en cambio, bajo la misma jefatura y el mismo funcionario que las tierras fiscales, el catastro y las obras públicas (capítulo \ref{cap:tierrafiscal}). \textbf{Los tres circuitos que este libro sigue para cada loteo ---el dominio, el plano y la línea de ribera--- dependen desde diciembre de 2025 de una sola oficina.} La estructura de esa jefatura, aprobada por la Decisión Administrativa 569/25 (23/12/2025, con vigencia al 15 de diciembre; B.O.\ Nº 22.100 del 30/12/2025), tiene seis secretarías ---Obras Públicas, Coordinación Interministerial, Interior, Prensa y Comunicación, Tierra y Bienes del Estado y Relaciones Internacionales---, y la Dirección General de Inmuebles depende directamente del jefe de Gabinete. \textbf{La Subsecretaría de Recursos Hídricos no depende de él sino de la Secretaría de Obras Públicas}: el organismo que determina las líneas de ribera quedó un escalón por debajo del que proyecta y contrata las obras, junto a sus direcciones de obras municipales y de infraestructura social y productiva. La Secretaría \textbf{del Interior}, por su parte, tiene a su cargo las relaciones con los municipios a través de una Subsecretaría de Asuntos Municipales. Su anexo nombra a los responsables del área hídrica, y dos ya aparecieron en este libro. El coordinador de la Subsecretaría es \textbf{David Le Favi}\index{Le Favi, David}, que en 2023, como director de Fiscalización y Control, firmó la nota a Fiscalía de Estado analizada en el capítulo \ref{cap:amparo}. El director general de Cuencas y Política Hídrica es \textbf{Jorge Genaro Torres}\index{Torres, Jorge G.}, con toda probabilidad el geólogo Jorge G.\ Torres de este capítulo: el autor de la presentación sobre línea de ribera y el integrante de las comisiones de 2014 y 2015. El director general de Concesiones Hídricas es Leandro Bruno.

Dos lecturas se siguen de ese cambio.

\textbf{La primera es de capacidad.} El capítulo \ref{cap:amparo} documenta que en la audiencia
de febrero de 2026 la propia Subsecretaría alegó falta de capacidad operativa. Que esa
alegación provenga de un organismo que acababa de perder rango ---y con él, presumiblemente,
presupuesto y planta--- no la vuelve más ni menos atendible, pero la ubica.

\textbf{La segunda es de plantilla.} El edicto lo emite la Subsecretaría y, sin embargo, ordena
presentar el recurso «ante la Secretaría de Recursos Hídricos». \textbf{El formulario no se
actualizó.} Es el mismo fenómeno que el error del «Dcto.\ 1898/02» que este capítulo documentó
entre 2013 y 2025: un texto que se arrastra de un edicto al siguiente sin que nadie lo lea.
\end{cautela}

\lamina{lam-da565-cabecera.jpg}{lam:da565} {Estructura del Ministerio de Producción y Minería desde el 15 de diciembre de 2025} {\textbf{Anexo I de la Decisión Administrativa 565/25, cabecera.} Bajo el ministro, cuatro secretarías: Ambiente y Desarrollo Sustentable, Industria y Comercio, Minería y Desarrollo Agropecuario. No hay unidad de recursos hídricos en ningún nivel del organigrama.} {Boletín Oficial de Salta Nº 22.100, 30/12/2025; anexo SA100052284, folio 3. Firmas del ministro José Ignacio Lupión y del coordinador administrativo Nicolás Demitrópulos. Escaneo en escala de grises, sin otras alteraciones.}

\textbf{Tres detalles menores del edicto de julio de 2025}, que conviene registrar porque cada uno cierra
algo que este capítulo dejó abierto.

\textbf{El número del decreto está bien citado} en este edicto, a diferencia del de abril del mismo año.

\textbf{El plazo se precisó.} Los diez días hábiles corren ahora «\textit{desde el día siguiente
de la última publicación}», y no «desde el día de la última publicación» como en 2013 y 2014. Es
un día de diferencia y es la clase de ajuste que sólo se hace cuando alguien lo discutió.

\textbf{Y la vía de recurso admite correo electrónico}: el edicto ofrece presentar el recurso en
Avenida Bolivia 4650 «\textit{o bien al mail}» de mesa de entradas de la Secretaría, que es una
casilla de Gmail. Este libro lo consigna sin ironía: \textbf{amplía el acceso de quien vive a
veinte kilómetros}, y es exactamente lo contrario de remitir las coordenadas a la sede. El mismo
organismo, en el mismo acto, acerca el recurso y aleja el dato.

\section{Cómo se ve un edicto de línea de ribera bien hecho}

El 7 y 8 de marzo de 2019, la misma Secretaría de Recursos Hídricos publicó un edicto sobre el río Bermejo, departamento Orán, expediente 0090034-289615/18. Sirve de patrón de comparación.

\begin{ficha}{Edicto del río Bermejo --- B.O.\ Nº 20.457, 7 y 8 de marzo de 2019}
Por \textbf{Resolución Nº 21 del 01/02/19 se aprobó la Comisión Técnica} que determinará la línea de ribera en ambas márgenes del río Bermejo, integrada por el \textbf{geólogo Enrique Chalabe} y el \textbf{ingeniero civil Marcelo Antonio Chalabe}. Se establece la zona: ambas márgenes del Bermejo en el área de la matrícula Nº 30.806 de Orán. \textbf{Método: geológico-geomorfológico.} Fundamento: Decreto Nº 1989/02, reglamentación del art.\ 126, Anexo I.

\textbf{``Los propietarios ribereños no podrán realizar ningún trabajo que implique la modificación actual de las márgenes en la zona objeto de la presente determinación.''}

Publicación por dos días. Los terceros interesados disponen de \textbf{quince días hábiles administrativos} desde la última publicación para presentar información, estudios, planos o cualquier dato que oriente la determinación, ante la sede de Avenida Bolivia 4.650. Firma: \textbf{Dra.\ Silvia Santamaría\index{Santamaría, Silvia F.}, Jefa de Programa Jurídico}.
\end{ficha}

\begin{cautela}
Este edicto obliga a corregir dos datos que el apartado anterior había tomado de un resumen intermedio, y a reformular la observación principal.

\textbf{Corrección primera:} la Comisión Técnica de este edicto tiene \textbf{dos} integrantes --- un geólogo y un ingeniero civil ---, cuando el artículo 3º del Anexo I del Decreto 1989/02 la conforma con \textbf{tres}: ingeniero hidráulico, geólogo y agrimensor. La diferencia no es un error de este libro sino un dato del expediente del Bermejo, y corresponde preguntar por ella, no darla por normal.

\textbf{Corrección segunda:} los quince días hábiles del edicto del Bermejo son para que terceros aporten estudios \textbf{durante} la determinación. Conviven con el plazo de diez días del artículo 10 del Anexo I, que corre para observar la línea propuesta y, cumplido sin oposición, la deja firme. Son dos plazos de dos momentos, no uno que desmiente al otro.

\textbf{Y la reformulación, que es más importante que las dos correcciones.} Lo que este edicto revela no es un plazo sino una \textbf{etapa}. Se publica \textbf{antes} de determinar: anuncia la comisión con nombre y apellido, declara el método que se empleará, delimita la zona, prohíbe cautelarmente a los ribereños modificar las márgenes mientras dure la operación, y abre quince días para que cualquiera aporte estudios. Eso es lo que el artículo 126 llama ``dar intervención en la operación a los interesados''.

El edicto de la Resolución 383/14 que este capítulo documenta pertenece al momento opuesto: comunica que la determinación \textbf{ya fue aprobada}, y remite las coordenadas a la sede. \textbf{No son el mismo acto y no deben compararse como si lo fueran.}

De donde surge una pregunta precisa: si el procedimiento incluye un edicto inicial como el del Bermejo, entonces para los arroyos El Durazno y Guaranguay \textbf{debió existir uno equivalente entre 2013 y 2014}, con los nombres de la comisión, el método elegido, la zona delimitada y la prohibición cautelar de modificar las márgenes. Si existió, hay que encontrarlo en el Boletín de esos años. Si no existió, la determinación se hizo sin la etapa participativa que la ley prevé. Las dos respuestas son informativas y ninguna está en poder de este libro.
\end{cautela}

\subsection*{Y el texto del Anexo I lo convierte en obligación}

El texto íntegro del Anexo I del Decreto 1.989/02 vuelve innecesaria la analogía con el edicto del Bermejo.

\begin{ficha}{Anexo I, artículos 2 y 4 --- texto literal}
\textbf{Art.\ 2.} «\textit{Partes interesadas. Serán partes interesadas todos los que se
consideren con un interés legítimo o en un derecho subjetivo y los organismos provinciales de
participación obligatoria.}»

\textbf{Art.\ 4.} «\textit{Notificación. La Comisión formada al efecto deberá notificar a
todas las partes interesadas que la Autoridad de Aplicación ha decidido realizar la
determinación de la línea de ribera correspondiente y además la Comisión \textbf{deberá
publicar en el Boletín Oficial por 2 (dos) días y en el medio gráfico de mayor difusión un
aviso de inicio de la determinación prevista}.}»
\end{ficha}

\textbf{El aviso de inicio no es una práctica: es un deber, con plazo y con soporte.} Y el
artículo 5 le da su razón de ser: las partes tienen quince días desde la notificación para
presentar «\textit{toda información, estudio, planos o cualquier otro dato}» que oriente a la
Comisión.

\begin{cautela}
De modo que la pregunta sobre el edicto inicial de El Durazno y Guaranguay tiene ahora consecuencia jurídica: \textbf{si no se publicó, la determinación se hizo omitiendo un acto que el Anexo I ordena en términos imperativos}. Este libro sigue sin saber cuál de las dos cosas ocurrió, pero ya no compara con otra cuenca: compara con el texto.
\end{cautela}

\subsection*{Y el aviso de inicio existe: está publicado en 2012}

Hasta acá la pregunta era si el aviso de inicio del artículo 4 se cumplía en la práctica. La
respuesta es que sí, y el ejemplo está en este mismo departamento, en el expediente contiguo.

\begin{ficha}{Boletín Oficial Nº 18.970, 13 de diciembre de 2012 --- O.P.\ Nº 100031653}
Referencia: \textbf{Expediente Nº 34-179.209/12}.

«La Secretaría de Recursos Hídricos hace saber que, por \textbf{Resolución Nº 322 del día
21/11/12 se aprobó la Comisión Técnica que determinará la línea de ribera} en ambas márgenes
del \textbf{Río La Caldera}, integrando dicha comisión los Sres.\ \textbf{Ing.\ Civil Felipe
Caro}, M.P.\ Nº 3024, el \textbf{Geólogo Marcelo Ramón Olañeta}, M.F.\ Nº 155, y el
\textbf{Topógrafo N.\ Félix Márquez}. Estableciéndose la determinación de
dicha línea, respecto del inmueble \textbf{Catastro Nº 3593, Fracción 17}, del Dpto.\ La
Caldera, como lugar donde la comisión técnica realizará dicha determinación,
\textbf{utilizándose el método de registro fotográfico y medición de caudales en base a datos
de precipitaciones. Decreto Nº 1989/02, art.\ 8 ap.\ 2 y 3}. Los propietarios ribereños en la
zona objeto de la presente determinación \textbf{no podrán realizar ningún trabajo que implique
la modificación actual de las márgenes}.»
\end{ficha}

\textbf{Eso es, punto por punto, lo que el artículo 4 manda publicar.} Y publica más de lo que
manda: los tres integrantes de la comisión con su matrícula profesional, el inmueble, la
restricción que pesa sobre los ribereños mientras dura el trámite, y ---lo que este libro venía
pidiendo desde el comienzo--- \textbf{el método elegido, con su fundamento normativo}.

\begin{cautela}
\textbf{Tres cosas se siguen de esto.}

\textbf{Primera: el método se declara, y no es el primero de la prelación.} El artículo 8
ordena preferir aforos con estadística no menor de cincuenta años; este edicto invoca
expresamente los \textbf{apartados 2 y 3} ---registro fotográfico y modelos---. No es una
irregularidad: el propio artículo prevé descender cuando los datos del primer escalón son
insuficientes o no existen. \textbf{Es la confirmación, en un acto oficial de 2012, de que el organismo tampoco contaba con una serie de aforos que alcanzara el umbral de cincuenta años}: las estaciones existieron y sus registros están publicados, pero la más larga no llega a ese mínimo.

\textbf{Segunda: el expediente 34-179.209/12 tiene sus dos publicaciones.} El aviso de inicio
en diciembre de 2012 y la determinación ---Resolución 6/14--- en enero de 2014, con coordenadas publicadas
---que resultan ser las de la línea del pueblo, según se explica en la cautela sobre las coordenadas---.

\textbf{Y tercera, y hay que formularla con precisión porque es más fina de lo que parece.}
El expediente \textbf{34-167.127/13} \textit{sí} tiene su acto de comisión publicado: es la
Resolución 395/13 de diciembre de 2013, que este capítulo ya consigna. \textbf{Pero ese acto
aprueba la comisión para el río La Caldera y el arroyo El Manzano}, que son los dos cursos que
la Resolución 185/14 determinó en julio de 2014 y publicó con ciento trece pares de
coordenadas.

\textbf{Los arroyos El Durazno y Guaranguay no están en él.} Y las búsquedas en el buscador
oficial de avisos ---por cada uno de los dos nombres y por el número de expediente, en septiembre de 2026--- no
devuelven ningún acto que apruebe una comisión para esos dos cursos.

\end{cautela}

\begin{cautela}
De modo que, dentro de un mismo expediente, hay dos determinaciones con trayectorias distintas.

\textbf{La del río La Caldera y El Manzano}: comisión aprobada y publicada en diciembre de
2013, determinación publicada en julio de 2014 con todas sus coordenadas.

\textbf{La de El Durazno y Guaranguay}: sin acto de comisión que los nombre, y determinación
publicada en diciembre de 2014 sin coordenadas.

\textbf{El mismo expediente, el mismo año, el mismo firmante, y los dos cursos que atraviesan
el loteo son los que quedaron afuera del circuito completo.} Este libro registra la
coincidencia y no le atribuye intención: pudo haber un acto de comisión no publicado, o
publicado con otro título, o los dos arroyos pudieron incorporarse al trámite después. Las tres
posibilidades se resuelven con el expediente, que sigue pedido.
\end{cautela}

\subsection*{Y el mismo arroyo tuvo sus coordenadas publicadas en 2012}

En la misma edición, dos columnas antes, hay un segundo edicto.

\begin{ficha}{Boletín Oficial Nº 18.970 --- O.P.\ Nº 400002609}
Referencia: \textbf{Expediente Nº 34-129.165/10}.

«La Secretaría de Recursos Hídricos hace saber que, por \textbf{Resolución Nº 328 del día
22/11/12}, se aprobó la determinación de la línea de ribera realizada por la Comisión Técnica
\textbf{en ambas márgenes del Arroyo Guaranguay}, integrando dicha comisión el \textbf{Dr.\ en
Geología Víctor Omar Nieva} y la \textbf{Ing.\ Civil Mariela Adriana Nieva}, a la latitud del
inmueble identificado con la \textbf{matrícula catastral Nº 4082} del Dpto.\ La Caldera,
\textbf{conforme las siguientes coordenadas:}»

Sigue la tabla, encabezada «\textit{Vinculación de Líneas de Ribera del Arroyo Guaranguay
---~Mat.\ 4.082}», con margen derecha e izquierda en \textbf{coordenadas Gauss Krüger}.

Firma el \textbf{Dr.\ Matías J.\ Brogin}, Asesor Legal. Importe: \textbf{\$120,00}. Publicado
el \textbf{12 y 13 de diciembre de 2012}.
\end{ficha}

\textbf{Es el mismo arroyo.} Guaranguay, ambas márgenes, departamento La Caldera. Dos años
antes de la Resolución 383/14, sobre otra matrícula y en otro expediente, \textbf{y con las
coordenadas publicadas en el edicto}.

Y lo firma la misma persona.

\subsection*{Y ocho meses antes, el aviso de inicio de esa misma determinación}

Falta una pieza, y estaba en abril del mismo año.

\begin{ficha}{Boletín Oficial Nº 18.816, 24 y 25 de abril de 2012 --- O.P.\ Nº 400001406}
Referencia: \textbf{Expediente Nº 34-21.387/12}.

«La Secretaría de Recursos Hídricos hace saber que por \textbf{Resolución Nº 72 del día
30/03/12, se aprobó la Comisión Técnica que determinará la línea de ribera en ambas márgenes
del Arroyo Guaranguay}, integrando dicha comisión los Sres.\ \textbf{Ingeniero Civil Mariela
Nieva} y \textbf{Dr.\ en Geología Víctor Omar Nieva}. Estableciéndose el inmueble
\textbf{Catastro 4082} del Depto.\ La Caldera, como el lugar donde la comisión técnica realizará
el estudio y determinación de la línea de ribera, \textbf{conforme art.\ 126 del Código de
Aguas, reglamentado por el Decreto Nº 1989/02}.

Se ordena la publicación del presente por el término de dos días \ldots\ Conforme C.A., art.\
309, \textbf{las personas que tengan interés legítimo podrán hacerlos valer, en el término de 15
días hábiles administrativos}, ante la Secretaría de Recursos Hídricos, sita en Av.\ Bolivia
Nº 4650.»

Firma el \textbf{Dr.\ Matías J.\ Brogin}, Asesor Legal. Importe: \textbf{\$120,00}.
\end{ficha}

\textbf{El procedimiento completo, sobre el arroyo Guaranguay, está publicado en 2012.}

\begin{ficha}{La secuencia del arroyo Guaranguay, matrícula 4082}
\textbf{Abril de 2012.} Aviso de inicio: se aprueba la comisión, se la nombra, se cita la norma
que funda el procedimiento y se abre el plazo de quince días para quien tenga interés legítimo.

\textbf{Diciembre de 2012.} Determinación: la misma comisión ---los dos Nieva--- concluye, y el
edicto publica las coordenadas de ambas márgenes.

\textbf{Marzo de 2013.} Rectificación de esa determinación, publicada por dos días.

\textbf{Todo firmado por el mismo asesor legal, y todo por ciento veinte pesos cada
publicación.}
\end{ficha}

\begin{cautela}
Puesto al lado de lo que este capítulo documenta para el expediente 34-167.127/13, el contraste
es completo y no requiere ningún comentario adicional.

\textbf{Sobre el arroyo Guaranguay, matrícula 4082}, entre 2012 y 2013: aviso de inicio con
comisión nombrada y plazo abierto, determinación con coordenadas publicadas, y rectificación
publicada.

\textbf{Sobre el arroyo Guaranguay, matrícula 4366}, en 2014: ninguna de las tres cosas, hasta donde alcanzó la búsqueda de este libro. Ni
aviso de inicio, ni coordenadas, ni comisión nombrada en el edicto.

\textbf{Mismo curso de agua. Mismo organismo. Mismo firmante. Dos años de diferencia.}

Este libro sigue sin afirmar por qué. Y ya no necesita afirmarlo: \textbf{la comparación es
interna, es del propio organismo consigo mismo, y está publicada en el Boletín Oficial de la
Provincia}. Cualquiera puede repetirla.
\end{cautela}

\begin{cautela}
\textbf{Con esto la serie de aranceles queda completa, y no deja nada en pie.}

\begin{itemize}[leftmargin=*,itemsep=0.2em,topsep=0.3em]
\item \textbf{Abril de 2012}, Guaranguay: aviso de inicio, \textbf{\$120}.
\item \textbf{Diciembre de 2012}, Guaranguay: coordenadas publicadas, \textbf{\$120}.
\item \textbf{Diciembre de 2013}, río La Caldera: once pares publicados, \textbf{\$180}.
\item \textbf{Julio de 2014}, río La Caldera y El Manzano: ciento trece pares publicados,
      \textbf{\$250}.
\item \textbf{Diciembre de 2014} (edicto de enero de 2015), El Durazno y Guaranguay: \textbf{ninguna coordenada}.
\end{itemize}

Ciento trece pares costaron doscientos cincuenta pesos. Los de diciembre de 2012 costaron ciento veinte.
\textbf{El costo marginal de publicar una coordenada es de unos pocos pesos}, y este libro ya
no puede sostener ---ni siquiera como hipótesis de prudencia--- que la omisión de diciembre de
2014 se explique por extensión, arancel o composición.

\textbf{Una objeción previsible, y la respuesta.} Son pesos de 2012, 2013 y 2014, y el arancel
del Boletín se actualizó en ese período como todo lo demás: comparar los importes entre sí tiene
el límite que tiene. \textbf{La comparación que importa es otra, y es interna a cada edicto}:
en julio de 2014, ciento trece pares de coordenadas se publicaron por doscientos cincuenta
pesos. Cinco meses después ---resolución de diciembre, edicto de enero de 2015---, no se publicó
ninguna.

\textbf{Lo que no puede reemplazarse es la explicación.} Este capítulo no sabe por qué se
eligió una forma u otra, y sigue sin afirmar que la elección haya sido deliberada. Lo que
ahora está establecido es que \textbf{hubo elección}: dos plantillas disponibles, las dos en
uso, ninguna norma que mande preferir una, y ningún acto que funde el apartamiento.
\end{cautela}

\subsection*{La serie empieza antes, y el índice del Boletín lo esconde}

Este capítulo fechó su serie entre octubre de 2013 y diciembre de 2015, y los edictos de 2012
la extendieron un año hacia atrás. El listado del buscador oficial la extiende cuatro años más,
y por una razón que conviene entender.

\begin{ficha}{Determinaciones anteriores, publicadas bajo otro tipo de aviso}
\textbf{Marzo de 2009.} «Secretaría de Recursos Hídricos --- Expte.\ Nº 34-9.791/08 ---
\textbf{determinac.\ línea de ribera cat.\ 2981}, La Caldera.»

\textbf{Septiembre de 2011.} «Expte.\ Nº 34-50.034 -- 4.034/11 --- \textbf{línea de ribera
inmueble matrícula 1491} del Dpto.\ La Caldera, \textbf{Finca El Buen Retiro}.»

\textbf{Octubre de 2011.} «\textbf{Línea de ribera catastro Nº 4027} --- Dpto.\ La Caldera ---
Expte.\ Nº 0050034-136060/2011-0.»

Los tres están publicados bajo el tipo de aviso \textbf{«concesiones de agua pública»}, no bajo
«determinaciones de línea de ribera».
\end{ficha}

\begin{cautela}
\textbf{Dos consecuencias, y la segunda obliga a matizar algo que este capítulo afirmó.}

\textbf{Primera: la serie es más larga.} Hay determinaciones de línea de ribera en este
departamento al menos desde 2009, sobre inmuebles que este libro no conocía ---catastro 2981,
matrícula 1491 de la finca El Buen Retiro, catastro 4027---.

\textbf{Segunda: el índice del Boletín no es una clasificación estable.} Un mismo tipo de acto
aparece, según el año, bajo «determinaciones de línea de ribera», «avisos administrativos» o
«concesiones de agua pública».

Eso obliga a matizar la afirmación sobre el aviso de inicio de El Durazno y Guaranguay.
\textbf{Las búsquedas de este libro fueron por texto} ---por el nombre de cada arroyo y por el
número de expediente---, y el texto atraviesa los tipos de aviso: así aparecieron los actos de
2012, que están bajo «avisos administrativos». De modo que la conclusión se sostiene.
\textbf{Pero un aviso mal indexado, con el nombre del arroyo escrito de otro modo o sin número
de expediente en el título, no aparecería}, y este libro no puede descartarlo.

\textbf{Por eso el apéndice \ref{cap:pedidos} pide una constancia escrita del organismo y no
una búsqueda más.} Una ausencia en un buscador es un indicio; una certificación firmada es un
documento.
\end{cautela}

\begin{ficha}{Listado del buscador de avisos --- «línea de ribera» y «La Caldera», consultado en septiembre de 2026}
Ciento siete avisos, de 2004 a enero de 2017. Además de los actos que este capítulo analiza, los títulos registran otros que no figuran en su tabla:
\begin{itemize}[leftmargin=*,nosep]
\item marzo y abril de 2004: «determinación línea de ribera», expediente 34-3.998/03;
\item septiembre de 2005: catastro 2961, sección B, Vaqueros; y diciembre de 2005, un aviso de la entonces Dirección General de Recursos Hídricos con la Resolución 683;
\item junio y julio de 2013, antes de la serie: \textbf{Resolución 155}, expediente 34-132.589/12, «arroyo Los Nogales --- matrícula 1905, sección A, fracción 17B», y expediente 34-214.573/12, «río Vaqueros, ambas márgenes --- matrícula catastral s/Nº, sección B --- localidad Vaqueros» (ambos transcriptos en el apartado sobre 2012 y 2013);
\item septiembre de 2011: \textbf{«línea de ribera Chacras del Gallinato», catastro 3709}, expediente 0050034-42754/2011. Abierto, es la comisión de las Resoluciones 396 y 418 de 2011; la determinación, por las Resoluciones 011 y 020 de 2012, se publicó en mayo de 2012 sin coordenadas (capítulo \ref{cap:loteo});
\item octubre y diciembre de 2011: otros dos expedientes sobre el catastro 4027 y el río Castellanos, además del ya transcripto;
\item junio y julio de 2015: \textbf{«Aguas del Norte --- Co.S.A.y S.A. --- río La Caldera --- Dpto.\ La Caldera»}. Abierto, el aviso resulta ser el de la Resolución 197/15, ya incluida en la tabla: comisión técnica (Torres y Cerezo) para ambas márgenes del río La Caldera en el catastro 4049, «método geológico-morfológico», edicto del 22 de junio de 2015. El texto no nombra a la empresa; sólo el índice del Boletín lo titula con ella;
\item septiembre de 2015 y marzo de 2016: expediente 34-125.733/15, sin curso ni lugar en el título. Abiertos, son la Resolución 273/15 (comisión, río Trancas, Finca Vaqueros), ya incorporada a la tabla, y la Resolución 14/16 que la cierra;
\item diciembre de 2016 y enero de 2017: expediente 34-202.169/15, «márgenes del río Vaqueros»: la Resolución 414/16, con coordenadas, cuya comisión (Res.\ 147/16) y ampliación (Res.\ 144/17) se transcriben más abajo.
\end{itemize}
\end{ficha}

\begin{cautela}
Esto acota tres afirmaciones de este capítulo. La serie de 2013 a 2015 es la que este libro reconstruyó con los textos leídos: el expediente 34-125.733/15, que el título no ubicaba, resultó ser una comisión del departamento y se sumó a la tabla, que con las Resoluciones 208/13 y 428/15, halladas después, pasa de veintidós a veinticinco actos; el aviso titulado con la empresa de agua ya estaba en ella con otro nombre. La serie de determinaciones del departamento no empieza en 2009 sino, al menos, en 2004. Y la línea de ribera pedida por el propio aparato estatal probablemente no aparece por primera vez en la reunión de 2017 sobre el Acueducto Norte. El índice del Boletín, que suele titular los avisos por quien los encarga y paga, asocia a la empresa provincial de agua con la comisión de 2015 sobre el río La Caldera en el catastro 4049; el decreto reglamentario pone los costos a cargo de quien solicita. Si el catastro 4049 es el de la planta de tratamiento, la determinación que la reunión de 2017 señalaba como paso previo ya tenía comisión designada dos años antes, con el mismo geólogo que luego participó de esa reunión, y lo que falta es el acto que la aprobó. Las búsquedas en el Boletín por el número de expediente y por el catastro, entre 2015 y 2019, devuelven sólo el aviso de la comisión: la determinación no se publicó con ninguno de esos dos datos en el título o el texto indexado. Los títulos no dicen si los demás actos son comisiones o determinaciones, ni si publicaron coordenadas; hasta leer los textos, este libro no los suma a sus recuentos. \pendiente{textos de los avisos de 2004 y 2005; solicitante de la Resolución 197/15, ubicación del catastro 4049 y resolución que aprobó esa determinación}

Y un límite del propio listado: ningún título es posterior a enero de 2017, aunque este capítulo transcribe actos de 2017, 2019 y 2025 sobre el mismo departamento. El caso de la Resolución 144/17 lo muestra: se publicó en junio de 2017 con el título «Secretaría de Recursos Hídricos --- Expte.\ Nº 34-202.169/15», sin el nombre del departamento, y por eso no aparece en una búsqueda por «La Caldera». El buscador sólo encuentra lo que el título dice.
\end{cautela}

\subsection*{Los tres edictos anteriores, leídos}

Abiertos, los tres contestan la pregunta que el capítulo tenía pendiente y agregan una que no
tenía.

\begin{ficha}{Boletín Oficial Nº 18.066, 10 y 11 de marzo de 2009}
«La Secretaría de Recursos Hídricos \ldots\ hace saber que \textbf{ha decidido autorizar la
determinación de la línea de ribera de la margen derecha con el Río Lesser} a la altura del
inmueble identificado como \textbf{Catastro Nº 2981 --- Fracción B8}, del Dpto.\ La Caldera, de
conformidad con el art.\ 126 de la Ley 7017/98 y su Decreto Reglamentario 1989/02.

Se hace saber que los interesados pueden formular oposiciones \ldots\ dentro de los
\textbf{15 días} a partir de la última publicación.»

Firma la \textbf{Dra.\ Silvia F.\ Santamaría}, Jefe Sub-Programa Coordinación y Capacitación.
Importe: \textbf{\$60,00}.
\end{ficha}

\begin{ficha}{Boletín Oficial Nº 18.687, 11 y 12 de octubre de 2011}
«\ldots\ ha decidido autorizar la determinación de la línea de ribera \textbf{del Río
Castellanos}, ambas márgenes, a la altura del \textbf{Catastro Nº 4027} \ldots

Se hace saber que \textbf{la Comisión Técnica designada} estará integrada por: \textbf{Ing.\
Civil Mariela Adriana Nieva y Dr.\ en Geología Víctor Omar Viera} \ldots\ los interesados pueden
formular oposiciones \ldots\ dentro de los 15 días.»

Firma la \textbf{Dra.\ Silvia F.\ Santamaría}. Importe: \textbf{\$120,00}.
\end{ficha}

\begin{ficha}{Boletín Oficial Nº 18.668, 7 y 8 de septiembre de 2011}
«Se ha determinado la línea de ribera del inmueble \textbf{matrícula 1491} del Dpto.\ La
Caldera, \textbf{finca El Buen Retiro}, cuyo polígono limita con el \textbf{arroyo La
Calderilla}. La Comisión Técnica fue integrada por el \textbf{Ing.\ Lucas Vidal} y
\textbf{Geólogo Héctor Saravia Navamuel} \ldots\ siendo la misma aprobada por \textbf{Resolución
Nº 417 de fecha 02/09/11}.»

Sigue la tabla: «\textit{Vinculación de Líneas de Ribera del arroyo La Calderilla --- Mat.\
1491}», \textbf{veinticuatro puntos} en coordenadas Gauss Krüger, doce por margen.

Firma el \textbf{Dr.\ Rafael Ángel Figueroa}\index{Figueroa, Rafael Ángel}, Jefe Programa Legal
y Técnico. Importe: \textbf{\$140,00}.
\end{ficha}

\textbf{Las dos plantillas están completas y funcionando en 2009 y 2011}, cuatro y dos años
antes del inicio de la serie que este capítulo reconstruyó: aviso de inicio con plazo de
oposición ---y, desde 2011, con la comisión nombrada--- y determinación con todas sus
coordenadas.

\begin{cautela}
\textbf{Tres observaciones.}

\textbf{Una firma que atraviesa diecisiete años.} La Dra.\ Silvia F.\ Santamaría firma el aviso
de 2009 como jefa de subprograma, el de 2011 igual, y en 2026 firma como Directora de la
Dirección Legal y Técnica los edictos de otros departamentos que este capítulo documenta.
\textbf{Se consigna el recorrido, no se le atribuye nada}: una funcionaria de carrera que
permanece es lo esperable en una administración, y este libro no tiene elemento alguno para
leerlo de otro modo.

\textbf{Una discrepancia de apellido.} El edicto de octubre de 2011 nombra al geólogo
\textbf{«Víctor Omar Viera»}; los de abril y diciembre de 2012 nombran a \textbf{«Víctor Omar
Nieva»}, junto a la misma ingeniera Mariela Nieva. Puede ser un error de tipeo en uno de los
dos, o dos personas distintas. Este libro no lo establece.

\textbf{Y un error dentro del propio acto.} El edicto de septiembre de 2011 dice que el polígono
limita con el \textbf{arroyo La Calderilla}, titula su tabla «\textit{Vinculación de Líneas de
Ribera del arroyo La Calderilla}» y, en el medio, afirma que la comisión «\textit{propone la
línea de ribera de ambas márgenes \textbf{del río Arenales}}».

\textbf{El río Arenales no pasa por este departamento}: es el de la ciudad de Salta. Es una
frase arrastrada de otro edicto, del mismo modo que el «Dcto.\ 1898/02» se arrastró durante más de una
década y que en 2026 un edicto de la Subsecretaría sigue mandando recurrir «ante la Secretaría».

\textbf{Este libro no infiere de ahí ninguna consecuencia sobre la validez del acto}, que no le
corresponde juzgar. Registra el hecho: en el instrumento que fija dónde termina la propiedad
privada y empieza el dominio público, sobre un inmueble concreto, figura el nombre de un río
equivocado.
\end{cautela}

\subsection*{Dos definiciones que faltaban}

El Anexo II del mismo decreto define lo que el artículo 7 manda medir, y una de las
definiciones es más exigente de lo que este libro suponía.

\begin{ficha}{Anexo II --- Conceptos y Denominaciones, punto 6}
\textbf{Crecidas ordinarias:} «\textit{son aquellas cuyos períodos de recurrencia son iguales o
menores a 50 años, es decir las crecidas que en período de 50 años se producen por lo menos dos
veces}».

\textbf{Crecidas extraordinarias:} recurrencia menor a 100 años.
\textbf{Crecidas excepcionales:} recurrencia mayor a 100 años.
\end{ficha}

\textbf{Eso explica el artículo 8.} La prelación que exige aforos con estadística no menor de
cincuenta años no es un capricho: es el período que la propia norma fija para definir qué
crecida es ordinaria. \textbf{Sin cincuenta años de registro no se puede aplicar el criterio
que la norma establece}, y por eso el artículo 8 admite descender a registro fotográfico y,
en su defecto, a modelos matemáticos.

Este capítulo mostró que ninguna estación de aforo opera hoy dentro de la cuenca y que la serie más larga no llega a los cincuenta años. \textbf{De eso se sigue que el primer escalón del artículo 8 no estaba disponible para ninguna de las quince determinaciones de 2013--2015, y que corresponde preguntar cuál de los otros dos se aplicó en cada una.}

\section{La apelación, y quién firma}

Del edicto de la 383/14, transcripto más arriba, importa ahora su pie de imprenta. \textbf{Se publicó los días 21 y 22 de enero de 2015 en los Boletines Oficiales Nº 19.466 y 19.467}, bajo orden de publicación 100045883 ---sin enlace «Ver anexo» en la ficha web---, y lleva fecha de firma del
20 de enero de 2015 --- \textbf{un mes y dos días después} de la resolución que publica.

Este texto permite precisar una observación sobre los plazos.

Los diez días hábiles \textbf{sí existen}: no son el plazo para participar en la determinación --- ése es el de quince días del edicto del Bermejo, en la etapa inicial --- sino el plazo para \textbf{apelar el acto ya dictado} ante el Tribunal de Aguas. Son dos plazos de dos etapas distintas, y la confusión fue de este libro, no de las fuentes.

Restituida la distinción, la observación se sostiene y es más precisa que antes. Un vecino de El Durazno que leyera el edicto en enero de 2015 disponía de diez días hábiles para apelar ante el Tribunal de Aguas una línea \textbf{cuyas coordenadas no estaban en el edicto}. Para saber si su parcela quedaba dentro o fuera tenía que ir a Avenida Bolivia 4650, pedir las coordenadas Posgar, hacerlas volcar sobre su plano y volver con un recurso. Diez días hábiles, desde La Caldera, para un trámite que empieza con un viaje a la capital.

El derecho de apelación existe y está expresamente indicado en el edicto, que es más de lo que muchos actos administrativos informan. Lo que este libro registra es la distancia entre el derecho y su ejercicio, que en este caso se mide en kilómetros y en un dato que se publicó por remisión.

\begin{cautela}
Dos precisiones documentales que el texto agrega.

\textbf{El decreto citado.} El edicto invoca el «Decreto 1898/02, artículo 10», con la transposición de dígitos que este capítulo ya documentó; el correcto es el 1989/02.

\textbf{La firma.} Este edicto lo suscribe el Asesor Legal Matías J.\ Brogin\index{Brogin, Matías J.}, no la funcionaria cuya trayectoria este capítulo sigue. La serie de firmas del área jurídica del organismo hídrico incluye entonces al menos cuatro nombres en el período examinado: Santamaría, Brogin, Rafael Ángel Figueroa y Sandra Mabel Siegrist, que firma la Resolución 360/16 y el edicto de Carante de 2014.
\end{cautela}

Y una serie que ya tiene cinco puntos. \textbf{Silvia Santamaría\index{Santamaría, Silvia F.}} firma avisos en 2009 y 2011 como jefa de subprograma; el edicto de la concesión de agua del Catastro 1700 de La Caldera en junio de 2014 como Abogada Jefe del Programa Fiscalización y Control; el edicto del Bermejo en marzo de 2019 como Jefa de Programa Jurídico; y figura como Directora Legal y Técnica de la Subsecretaría en la nómina de diciembre de 2025, cargo con el que firma en 2026. Diecisiete años en el área jurídica del organismo hídrico provincial, con ascensos sucesivos.

Se consigna, como en el capítulo \ref{cap:politica}, sin ninguna connotación: es el mapa de dónde está la memoria institucional, no un señalamiento.
